%% Volume I, Chapter 15: Vector Spaces
\chapter{Vector Spaces}
\label{ch:vector-spaces}

\abstract{Polynomials, continuous functions, sequences, matrices and the solutions of a linear differential equation have little in common at first sight. Yet all of them can be added and multiplied by numbers, and they obey the same eight rules as the vectors of $\mathbb R^n$. Taking these rules as axioms pays off at once: a theorem proved once holds for all of them. A single number, the dimension, counts the degrees of freedom of such a space, and four subspaces attached to a matrix answer every question about a linear system: whether it can be solved, for which right-hand sides, and how many solutions it has. They also explain why the rank, so easy to define, is so delicate to compute.}

\Cref{ch:vectors-matrices} built the arithmetic of vectors and matrices, and \cref{ch:linear-systems} solved linear systems by elimination and attached to every matrix a number, its rank, the number of pivots of its reduced row echelon form. Both chapters spoke only of lists of numbers, the vectors of $\mathbb R^n$, and both left questions open that need a more general language. \Cref{ch:vectors-matrices} observed that a third vector enlarges the span of two vectors in $\mathbb R^3$ to the whole space unless it lies in their plane, and asked what the dimension of a span is; \cref{ch:linear-systems} stated without a complete proof that a matrix and its transpose have the same rank, and anticipated that the special solutions of a homogeneous system form a basis of its null space. This chapter supplies the language in which these statements become precise: vector spaces, subspaces, linear independence, bases and dimension. With it, the rank becomes a dimension, and the equality of the ranks of $A$ and $A^{\mathsf T}$ becomes a theorem about four subspaces attached to $A$.

The guiding idea is abstraction for the sake of economy. The eight rules of \cref{prop:vector-rules} are all that most arguments about vectors in $\mathbb R^n$ use, and many other collections of objects obey them: the polynomials, the continuous functions on an interval, the infinite sequences, the matrices of a given shape, the solutions of the differential equation $f''=-f$ of \cref{ch:transcendental}. By taking the eight rules as axioms, we prove each theorem once and apply it to all of these at the same time. The price is a little more care in the definitions; the reward is that the reader who has understood $\mathbb R^2$ and $\mathbb R^3$ can see all the other spaces with the same eyes. The second theme is counting. The central theorem of the chapter says that all bases of a space have the same number of elements, and the proof reduces this to the fact of \cref{ch:linear-systems} that a homogeneous system with more unknowns than equations has a nonzero solution. The resulting number, the dimension, counts the degrees of freedom of a space, and most of the chapter consists of computing it.

Throughout the chapter $\mathbb F$ denotes either the field $\mathbb R$ of real numbers or the field $\mathbb C$ of complex numbers, and $\mathbb F^n$ and $\mathbb F^{m\times n}$ are the column vectors and the matrices with entries in $\mathbb F$. Everything holds for both fields with the same proofs, since only the rules of arithmetic of a field are used, and the results of \cref{ch:vectors-matrices,ch:linear-systems} that we quote hold for complex entries as well, as those chapters explained (\cref{sec:complex-matrices} and the introduction to \cref{ch:linear-systems}). The note at the end of \cref{sec:vector-spaces} explains what changes over other fields. Linear maps between abstract vector spaces, their kernels and images and the rank--nullity theorem in its general form are the subject of Chapter~16; here we need linear maps only in the form of matrices, $\mathbf x\mapsto A\mathbf x$, and of the coordinate maps of \cref{sec:coordinates}. Inner products appear only where they cost nothing; orthogonality of subspaces is the subject of Chapter~18.

The chapter is organised as follows. \Cref{sec:vector-spaces} states the axioms, derives their first consequences and collects examples, and \cref{sec:subspaces} studies subspaces, their intersections and sums, and the shifted subspaces, or flats, that are the solution sets of linear systems. \Cref{sec:independence} introduces spans and linear independence in general vector spaces, including spaces of functions, and \cref{sec:basis-dimension} proves that all bases of a space have the same number of elements and computes the dimensions of the examples. \Cref{sec:coordinates} shows that a basis turns every finite-dimensional space into a copy of $\mathbb F^n$ and how coordinates change with the basis. \Cref{sec:four-subspaces} attaches four subspaces to a matrix, reads off bases of all of them from the reduced row echelon form and proves that a matrix and its transpose have the same rank; \cref{sec:rank-revisited} collects the many faces of the rank. \Cref{sec:fredholm} characterises the right-hand sides for which a system is solvable and applies the result to the incidence matrices of graphs, and \cref{sec:computing-rank} turns the theory into an algorithm and shows by experiment why the rank of a matrix stored in floating-point arithmetic is a delicate quantity. Unlike most chapters of this volume, every result of the chapter is proved.

\section{Vector Spaces}
\label{sec:vector-spaces}

\subsection{The Axioms}

The vectors of $\mathbb R^n$ can be added and multiplied by real numbers, and \cref{prop:vector-rules} listed eight rules that these operations satisfy. A vector space is any set with two such operations that satisfy the same eight rules.

\begin{definition}[Vector space]
\label{def:vector-space}
A \term{vector space} over $\mathbb F$ is a set $V$ together with an element $\mathbf 0\in V$ and two operations, an \term{addition} $V\times V\to V$, $(\mathbf u,\mathbf v)\mapsto\mathbf u+\mathbf v$, and a \term{multiplication by scalars} $\mathbb F\times V\to V$, $(t,\mathbf v)\mapsto t\mathbf v$, such that for all $\mathbf u,\mathbf v,\mathbf w\in V$ and all $s,t\in\mathbb F$:
\begin{enumerate}
\item $(\mathbf u+\mathbf v)+\mathbf w=\mathbf u+(\mathbf v+\mathbf w)$;
\item $\mathbf u+\mathbf v=\mathbf v+\mathbf u$;
\item $\mathbf v+\mathbf 0=\mathbf v$;
\item there is an element $\mathbf w\in V$ with $\mathbf v+\mathbf w=\mathbf 0$;
\item $s(t\mathbf v)=(st)\mathbf v$;
\item $1\mathbf v=\mathbf v$;
\item $t(\mathbf u+\mathbf v)=t\mathbf u+t\mathbf v$;
\item $(s+t)\mathbf v=s\mathbf v+t\mathbf v$.
\end{enumerate}
The elements of $V$ are called \term{vectors}, the elements of $\mathbb F$ \term{scalars}, and $\mathbf 0$ is the \term{zero vector}. A vector space over $\mathbb R$ is a \term{real vector space}, one over $\mathbb C$ a \term{complex vector space}.
\end{definition}

The axioms are the rules (i) to (viii) of \cref{prop:vector-rules}, in the same order, with one change: (iv) asserts the existence of a negative instead of naming it, since in an abstract space there is no componentwise formula for $-\mathbf v$. Two requirements are hidden in the word \emph{operation}: the sum of two elements of $V$ and the multiple of an element of $V$ must again lie in $V$. A set that is not closed under the operations is not a vector space, however well the rules hold where they make sense. We write the elements of an abstract vector space in bold, as we do for the vectors of $\mathbb F^n$, when we think of them as vectors; when they are functions, polynomials or matrices, we keep their usual names $f$, $p$ and $A$.

The familiar rules of computation follow from the axioms. \Cref{sec:vectors-in-rn} derived some of them from the eight rules; we collect them here, with proofs that use nothing but the axioms, so that they hold in every vector space.

\begin{proposition}[Consequences of the axioms]
\label{prop:vector-space-rules}
Let $V$ be a vector space over $\mathbb F$, $\mathbf u,\mathbf v,\mathbf w\in V$ and $t\in\mathbb F$.
\begin{enumerate}
\item The zero vector is unique: if $\mathbf v+\mathbf z=\mathbf v$ for all $\mathbf v$, then $\mathbf z=\mathbf 0$.
\item The negative is unique: there is exactly one $\mathbf w$ with $\mathbf v+\mathbf w=\mathbf 0$. It is denoted $-\mathbf v$, and $\mathbf u-\mathbf v\coloneqq\mathbf u+(-\mathbf v)$.
\item Cancellation: if $\mathbf u+\mathbf w=\mathbf v+\mathbf w$, then $\mathbf u=\mathbf v$.
\item $0\mathbf v=\mathbf 0$ and $t\mathbf 0=\mathbf 0$.
\item $(-1)\mathbf v=-\mathbf v$, and $-(-\mathbf v)=\mathbf v$.
\item If $t\mathbf v=\mathbf 0$, then $t=0$ or $\mathbf v=\mathbf 0$.
\end{enumerate}
\end{proposition}

\begin{proof}
(i) Taking $\mathbf v=\mathbf 0$ gives $\mathbf 0+\mathbf z=\mathbf 0$, and by (ii) and (iii) of \cref{def:vector-space}, $\mathbf 0+\mathbf z=\mathbf z+\mathbf 0=\mathbf z$; so $\mathbf z=\mathbf 0$.

(ii) If $\mathbf v+\mathbf w=\mathbf 0$ and $\mathbf v+\mathbf w'=\mathbf 0$, then, by the axioms (iii), (ii) and (i),
\[
\mathbf w=\mathbf w+\mathbf 0=\mathbf w+(\mathbf v+\mathbf w')=(\mathbf w+\mathbf v)+\mathbf w'=(\mathbf v+\mathbf w)+\mathbf w'=\mathbf 0+\mathbf w'=\mathbf w'.
\]

(iii) Add $-\mathbf w$ to both sides and use associativity: $\mathbf u=\mathbf u+(\mathbf w-\mathbf w)=(\mathbf u+\mathbf w)-\mathbf w=(\mathbf v+\mathbf w)-\mathbf w=\mathbf v$.

(iv) By axiom (viii), $0\mathbf v+0\mathbf v=(0+0)\mathbf v=0\mathbf v=0\mathbf v+\mathbf 0$, and cancellation gives $0\mathbf v=\mathbf 0$. Similarly $t\mathbf 0+t\mathbf 0=t(\mathbf 0+\mathbf 0)=t\mathbf 0=t\mathbf 0+\mathbf 0$ by axiom (vii), so $t\mathbf 0=\mathbf 0$.

(v) By the axioms (vi) and (viii) and part (iv), $\mathbf v+(-1)\mathbf v=1\mathbf v+(-1)\mathbf v=(1-1)\mathbf v=0\mathbf v=\mathbf 0$, so $(-1)\mathbf v$ is the negative of $\mathbf v$ by (ii). Since $-\mathbf v+\mathbf v=\mathbf 0$, $\mathbf v$ is the negative of $-\mathbf v$.

(vi) If $t\ne0$, then $\mathbf v=1\mathbf v=(t^{-1}t)\mathbf v=t^{-1}(t\mathbf v)=t^{-1}\mathbf 0=\mathbf 0$, by the axioms (vi) and (v) and part (iv).
\end{proof}

Part (vi) is the first place where it matters that the scalars form a field, in which every nonzero number has an inverse. These derivations are not difficult, and after this proposition we compute in a vector space exactly as in $\mathbb R^n$, without further reference to the axioms. Sums of several vectors need no parentheses, by associativity, and their order does not matter, by commutativity, for the reasons explained in \cref{sec:reals} for sums of numbers.

\subsection{Examples}

The point of the definition is the variety of its examples. In each of the following, the operations are the obvious ones, and checking the axioms amounts to the observation that they hold in each component, at each point or for each coefficient, because they hold for numbers.

\emph{Column vectors and matrices.} $\mathbb F^n$, with addition and multiplication by scalars component by component, is a vector space by \cref{prop:vector-rules}, whose proof works verbatim for complex numbers. The same is true of the set $\mathbb F^{m\times n}$ of $m\times n$ matrices, with the entrywise operations of \cref{sec:matrices-linear-maps}: as far as addition and multiplication by scalars are concerned, a matrix is a list of $mn$ numbers written in a rectangle, and the shape plays no role until matrices are multiplied with each other. The space $\mathbb C^n$ can also be regarded as a real vector space, by allowing only real scalars; we shall see that its dimension is then $2n$ rather than $n$, so the field of scalars is part of the structure.

\emph{Functions.} For a nonempty set $S$, the set of all functions $f\colon S\to\mathbb F$ is a vector space with the pointwise operations
\[
(f+g)(x)\coloneqq f(x)+g(x),\qquad(tf)(x)\coloneqq t\,f(x)\qquad(x\in S),
\]
whose zero vector is the function that is $0$ everywhere and in which $-f$ is the function $x\mapsto-f(x)$. The axioms hold because they hold at every point $x$; for instance, $((f+g)+h)(x)=(f(x)+g(x))+h(x)=f(x)+(g(x)+h(x))=(f+(g+h))(x)$ for every $x$, and two functions that agree at every point are equal. This single example contains many others. For $S=\{1,\ldots,n\}$, a function $i\mapsto x_i$ is a list of $n$ numbers, and we recover $\mathbb F^n$. For $S=\mathbb N$, a function $k\mapsto a_k$ is a sequence $(a_0,a_1,a_2,\ldots)$, and the sequences form a vector space, denoted $\mathbb F^{\mathbb N}$, with termwise addition and multiplication by scalars. For $S$ an interval, we obtain the space of all real functions on the interval, most of which are too wild to be of any interest; the interesting spaces are inside it, as the next examples show.

\emph{Continuous and differentiable functions.} The continuous real functions on an interval $[a,b]$ form the space $C[a,b]$, and the functions with $k$ continuous derivatives on an interval $I$ form the space $C^k(I)$. The operations are the pointwise operations of the previous example, and the only new point is that they do not lead out of the set: sums and multiples of continuous functions are continuous (\cref{ch:limits}), and sums and multiples of differentiable functions are differentiable, with $(f+g)'=f'+g'$ and $(tf)'=tf'$ (\cref{ch:differentiation}). The axioms then hold because they hold for all functions. \Cref{sec:subspaces} turns this argument into a general principle.

\emph{Polynomials.} A \term{polynomial} with coefficients in $\mathbb F$ is a function $p\colon\mathbb F\to\mathbb F$ of the form $p(x)=c_0+c_1x+\cdots+c_nx^n$. The coefficients are determined by the function, since a nonzero polynomial has only finitely many roots (\cref{cor:factor}), so that $p-q=0$ forces all coefficients of $p-q$ to vanish; the \term{degree} of $p\ne0$ is therefore well defined, as the largest $k$ with $c_k\ne0$. We denote by $\mathcal P$ the vector space of all polynomials and by $\mathcal P_n$ the polynomials of degree at most $n$, including the zero polynomial; sums and multiples of polynomials of degree at most $n$ have degree at most $n$, with the coefficients added and multiplied. Whether the coefficients are real or complex will be clear from the context or will not matter. The calligraphic letter distinguishes these spaces from the permutation matrices $P$ of \cref{ch:linear-systems}; with an argument, $\mathcal P(A)$ remains the power set of \cref{ch:sets}.

\emph{Solutions of linear differential equations.} Let $\omega>0$. If $f$ and $g$ are twice differentiable and satisfy $f''=-\omega^2f$ and $g''=-\omega^2g$, then $(f+g)''=f''+g''=-\omega^2(f+g)$ and $(tf)''=-\omega^2(tf)$; so the solutions of the equation form a vector space, and the same holds for every equation of the form $a_kf^{(k)}+\cdots+a_1f'+a_0f=0$ with constant coefficients. The equation $f''=-f$ is the one with which \cref{ch:transcendental} defined the sine and the cosine, and \cref{lem:harmonic-unique} showed that its solutions are exactly the functions $f(x)=a\cos x+b\sin x$. Applied to $g(t)\coloneqq f(t/\omega)$, which satisfies $g''=-g$, the lemma shows that the solutions of $f''=-\omega^2f$ are the functions $a\cos\omega x+b\sin\omega x$: the space of solutions is described by two numbers, which \cref{sec:basis-dimension} will express by saying that it has dimension~$2$. Similarly, the real sequences that satisfy the recurrence $a_{k+2}=a_{k+1}+a_k$ of the Fibonacci numbers form a vector space of sequences, and each of them is determined by its first two terms. Part~IV of this volume develops this observation for linear differential equations and recurrences of every order.

\emph{Positive numbers with unusual operations.} On the set $\mathbb R_{>0}$ of positive real numbers, define the addition $x\oplus y\coloneqq xy$ and the multiplication by real scalars $t\odot x\coloneqq x^t$. This is a real vector space: its zero vector is the number $1$, the negative of $x$ is $1/x$, and the axioms are the laws of exponents of \cref{thm:exponent-laws}; for instance $(s+t)\odot x=x^{s+t}=x^sx^t=(s\odot x)\oplus(t\odot x)$ and $t\odot(x\oplus y)=(xy)^t=x^ty^t$. The example shows that the axioms are about the operations, not about arrows or lists; \cref{sec:coordinates} shows that the logarithm turns this space into a copy of $\mathbb R$.

\emph{Sets that are not vector spaces.} The first quadrant $\{(x,y)\in\mathbb R^2: x,y\ge0\}$ is closed under addition but not under multiplication by $-1$. A line in $\mathbb R^2$ that does not pass through the origin, such as $x+y=1$, does not contain the zero vector, and the sum of two of its points, such as $(1,0)+(0,1)=(1,1)$, leaves it. The polynomials of degree exactly $2$ do not form a vector space, since $x^2+(1-x^2)=1$ has degree $0$. The singular $2\times2$ matrices are closed under multiplication by scalars but not under addition: $\operatorname{diag}(1,0)$ and $\operatorname{diag}(0,1)$ are singular, and their sum is the identity. The solutions of the inhomogeneous equation $f''=-f+1$ are not closed under addition either, since the sum of two solutions satisfies $f''=-f+2$. In each case one of the two operations leads out of the set, and in each case the set does not contain the zero vector; \cref{sec:subspaces} shows that these two defects are closely related.

\Cref{tab:vector-spaces} collects the examples, together with their dimensions, the number that \cref{sec:basis-dimension} attaches to every vector space and computes for these. Four of the spaces are infinite-dimensional, which means that no finite list of their elements generates all of them; in this volume such spaces serve mainly as surroundings for the finite-dimensional spaces inside them, and Volume~II studies them in their own right.

\begin{table}
\sdfloat{%
\begin{sdtable}{colspec={lll}, row{1}={font=\small\bfseries}}
\toprule
space & elements & dimension \\
\midrule
$\mathbb F^n$ & columns of $n$ scalars & $n$ \\
$\mathbb F^{m\times n}$ & $m\times n$ matrices & $mn$ \\
symmetric $n\times n$ matrices & $A^{\mathsf T}=A$ & $n(n+1)/2$ \\
$\mathcal P_n$ & polynomials of degree at most $n$ & $n+1$ \\
solutions of $f''=-\omega^2f$ & $a\cos\omega x+b\sin\omega x$ & $2$ \\
Fibonacci-type sequences & $a_{k+2}=a_{k+1}+a_k$ & $2$ \\
$\mathcal N(A)$, $A\in\mathbb F^{m\times n}$ & solutions of $A\mathbf x=\mathbf 0$ & $n-\operatorname{rank}A$ \\
$\mathbb R_{>0}$ with $\oplus$, $\odot$ & positive numbers & $1$ \\
$\mathcal P$ & all polynomials & infinite \\
$C[a,b]$ & continuous functions & infinite \\
$\mathbb F^{\mathbb N}$ & sequences $(a_0,a_1,\ldots)$ & infinite \\
\bottomrule
\end{sdtable}}{Examples of vector spaces and their dimensions. The dimensions are computed in \cref{sec:basis-dimension}, that of $\mathcal N(A)$ in \cref{sec:four-subspaces}.\label{tab:vector-spaces}}
\end{table}

\begin{note}{Other fields of scalars}
The algebraic part of this chapter uses no properties of $\mathbb R$ or $\mathbb C$ beyond the axioms of a field (\cref{def:ring}): the definitions and theorems that do not involve functions of a real variable, inner products or norms hold verbatim for vector spaces over an arbitrary field $\mathbb F$, with one exception, the dimension of the space of antisymmetric matrices, which is pointed out where it occurs. Over the field $\mathbb F_2=\mathbb Z/2\mathbb Z$ of \cref{cor:zp-field}, a vector space of dimension $n$ has exactly $2^n$ elements, and the subspaces of $\mathbb F_2^n$ are the \emph{linear codes} of coding theory: a message of $k$ bits is encoded as a vector of a $k$-dimensional subspace of $\mathbb F_2^n$, chosen so that any two of its vectors differ in many positions, and errors in transmission are detected and corrected by the parity checks that define the subspace, a homogeneous linear system over $\mathbb F_2$~\cite{macwilliams1977}. The real numbers are themselves a vector space over the field $\mathbb Q$ of rational numbers, of infinite dimension, an example that the note in \cref{sec:basis-dimension} takes up. When the scalars are only required to form a commutative ring, such as $\mathbb Z$, the structure is called a \emph{module}, and much of the theory breaks down, because one cannot divide: in the $\mathbb Z$-module $\mathbb Z$ the elements $2$ and $3$ satisfy the relation $3\cdot2-2\cdot3=0$, yet neither is a multiple of the other, the module $\mathbb Z/2\mathbb Z$ has no basis at all, and part~(vi) of \cref{prop:vector-space-rules} fails, since $2\cdot[1]=[0]$ in $\mathbb Z/2\mathbb Z$. Integer matrices are reduced by the Hermite and Smith normal forms mentioned in \cref{ch:linear-systems}.
\end{note}

\section{Subspaces}
\label{sec:subspaces}

\subsection{Definition and Test}

Most vector spaces of interest are found inside larger ones, as the continuous functions are found among all functions and the solutions of $A\mathbf x=\mathbf 0$ among all vectors of $\mathbb F^n$.

\begin{definition}[Subspace]
\label{def:subspace}
A subset $U$ of a vector space $V$ is a \term{subspace} of $V$ if it is itself a vector space with the addition and the multiplication by scalars of $V$.
\end{definition}

To check that a subset is a subspace, one need not verify the eight axioms one by one. The identities among them hold in $U$ because they hold in $V$, and only the closure of $U$ under the operations needs checking.

\begin{theorem}[Subspace test]
\label{thm:subspace-test}
A subset $U$ of a vector space $V$ is a subspace if and only if
\begin{enumerate}
\item $\mathbf 0\in U$;
\item $\mathbf u+\mathbf w\in U$ for all $\mathbf u,\mathbf w\in U$;
\item $t\mathbf u\in U$ for all $\mathbf u\in U$ and $t\in\mathbb F$.
\end{enumerate}
Condition (i) may be replaced by the condition that $U$ is not empty.
\end{theorem}

\begin{proof}
Suppose that $U$ satisfies (i) to (iii). By (ii) and (iii), the addition and the multiplication by scalars of $V$ restrict to operations on $U$. The axioms (i), (ii), (v), (vi), (vii) and (viii) of \cref{def:vector-space} are identities that hold for all vectors of $V$, so in particular for those of $U$. Axiom (iii) holds in $U$ with the zero vector $\mathbf 0\in U$, and for axiom (iv), the negative $-\mathbf u=(-1)\mathbf u$ of $\mathbf u\in U$ lies in $U$ by (iii) and \cref{prop:vector-space-rules}\,(v). So $U$ is a vector space.

Conversely, let $U$ be a subspace, with zero vector $\mathbf z$. It is not empty, it is closed under the operations, which is (ii) and (iii), and $\mathbf z+\mathbf z=\mathbf z=\mathbf z+\mathbf 0$ in $V$, so that $\mathbf z=\mathbf 0$ by cancellation in $V$; this is (i). Finally, if $U$ is not empty and satisfies (iii), it contains $0\mathbf u=\mathbf 0$ for any of its elements $\mathbf u$, by \cref{prop:vector-space-rules}\,(iv).
\end{proof}

The test explains the remark at the end of \cref{sec:vector-spaces}: a subset that is closed under multiplication by scalars and not empty must contain $\mathbf 0$, so a set that misses the zero vector cannot be a subspace, and checking for $\mathbf 0$ is the quickest way to rule out a candidate. Every vector space $V$ has the two \term{trivial} subspaces $\{\mathbf 0\}$ and $V$. A subspace of a subspace of $V$ is a subspace of $V$.

\subsection{Examples}

The test makes the examples of \cref{sec:vector-spaces} easy to justify, and it produces many more.

\emph{Lines and planes through the origin.} For $\mathbf v\in V$, the set $\{t\mathbf v: t\in\mathbb F\}$ is a subspace: it contains $0\mathbf v=\mathbf 0$, and $s\mathbf v+t\mathbf v=(s+t)\mathbf v$ and $s(t\mathbf v)=(st)\mathbf v$ lie in it. For $\mathbf v\ne\mathbf 0$ in $\mathbb R^2$ or $\mathbb R^3$ it is the line through the origin in the direction $\mathbf v$. More generally, every span $\operatorname{span}\{\mathbf v_1,\ldots,\mathbf v_k\}$ in $\mathbb F^n$ is a subspace, by \cref{prop:span}, and so is every plane $\{\mathbf x\in\mathbb R^3: \langle\mathbf a,\mathbf x\rangle=0\}$ through the origin.

\emph{The null space and the column space.} For $A\in\mathbb F^{m\times n}$, the null space $\mathcal N(A)=\{\mathbf x\in\mathbb F^n: A\mathbf x=\mathbf 0\}$ of \cref{sec:systems} is a subspace of $\mathbb F^n$; this is exactly the content of \cref{thm:solution-structure}\,(i). The \term{column space}
\[
\mathcal C(A)\coloneqq\{A\mathbf x: \mathbf x\in\mathbb F^n\}=\operatorname{span}\{\mathbf a_1,\ldots,\mathbf a_n\},
\]
the set of all combinations $x_1\mathbf a_1+\cdots+x_n\mathbf a_n$ of the columns of $A$, is a subspace of $\mathbb F^m$, again by \cref{prop:span}. It is the image of the map $\mathbf x\mapsto A\mathbf x$, and by the column picture of \cref{sec:systems}, \emph{the system $A\mathbf x=\mathbf b$ is consistent if and only if $\mathbf b\in\mathcal C(A)$}. The two spaces live in different places: $\mathcal N(A)$ consists of inputs of the map, vectors with $n$ components, and $\mathcal C(A)$ of outputs, vectors with $m$ components. \Cref{sec:four-subspaces} adds two further subspaces to these.

\emph{Spaces of matrices.} In $\mathbb F^{n\times n}$, the symmetric matrices form a subspace, since $(A+B)^{\mathsf T}=A^{\mathsf T}+B^{\mathsf T}$ and $(tA)^{\mathsf T}=tA^{\mathsf T}$ (\cref{prop:transpose-rules}); so do the antisymmetric matrices, the diagonal matrices, the upper triangular matrices and the matrices with trace zero, the last because the trace is additive and homogeneous (\cref{prop:trace}). The invertible matrices do not form a subspace, since they do not contain the zero matrix, and neither do the singular ones, as \cref{sec:vector-spaces} showed.

\emph{Spaces of functions.} In the space of all real functions on an interval $I$, the continuous functions form a subspace, and $C^k(I)$ is a subspace of $C(I)$ for every $k\ge1$; this is the argument of \cref{sec:vector-spaces}, now in the form of the test. The polynomials, restricted to $I$, form a subspace of $C^\infty(I)$, and we obtain a chain of subspaces
\[
\mathcal P_0\subseteq\mathcal P_1\subseteq\mathcal P_2\subseteq\cdots\subseteq\mathcal P\subseteq C^\infty(I)\subseteq\cdots\subseteq C^1(I)\subseteq C(I).
\]
The even functions, $f(-x)=f(x)$, and the odd functions, $f(-x)=-f(x)$, on $\mathbb R$ form subspaces, and so do the functions with $f(0)=0$, but not those with $f(0)=1$. The solutions of a homogeneous linear differential equation form a subspace of $C^2(\mathbb R)$, as shown in \cref{sec:vector-spaces}, and in the sequence space $\mathbb F^{\mathbb N}$, the convergent sequences, the sequences with only finitely many nonzero terms and the sequences that satisfy a fixed linear recurrence $a_{k+2}=pa_{k+1}+qa_k$ are subspaces, the first by the rules for limits of \cref{ch:limits}.

\subsection{Intersections, Unions and Sums}

New subspaces can be made from old ones by intersection and by addition, but not by union.

\begin{proposition}[Intersections and unions]
\label{prop:intersection}
Let $V$ be a vector space.
\begin{enumerate}
\item The intersection of any family of subspaces of $V$ is a subspace of $V$.
\item The union of two subspaces $U$ and $W$ is a subspace if and only if $U\subseteq W$ or $W\subseteq U$.
\end{enumerate}
\end{proposition}

\begin{proof}
(i) Every subspace of the family contains $\mathbf 0$, and if $\mathbf u$ and $\mathbf w$ lie in every subspace of the family, so do $\mathbf u+\mathbf w$ and $t\mathbf u$. (ii) If one subspace contains the other, the union is the larger one. Otherwise there are vectors $\mathbf u\in U$ with $\mathbf u\notin W$ and $\mathbf w\in W$ with $\mathbf w\notin U$. If $U\cup W$ were a subspace, it would contain $\mathbf u+\mathbf w$, which then lies in $U$ or in $W$. If $\mathbf u+\mathbf w\in U$, then $\mathbf w=(\mathbf u+\mathbf w)-\mathbf u\in U$, and if $\mathbf u+\mathbf w\in W$, then $\mathbf u=(\mathbf u+\mathbf w)-\mathbf w\in W$; both contradict the choice of $\mathbf u$ and $\mathbf w$.
\end{proof}

The union of the two coordinate axes of $\mathbb R^2$ is the standard example: it contains $(1,0)$ and $(0,1)$ but not their sum. The intersection of the null spaces of two matrices $A\in\mathbb F^{m\times n}$ and $B\in\mathbb F^{p\times n}$ is the null space of the matrix obtained by stacking $A$ on top of $B$, the solution set of both systems together; two planes through the origin of $\mathbb R^3$ that do not coincide meet in a line through the origin, as panel~B of \cref{fig:subspaces} shows. The smallest subspace that contains two given subspaces is not their union but their sum.

\begin{definition}[Sum and direct sum]
\label{def:sum}
For subspaces $U$ and $W$ of $V$, the \term{sum} is
\[
U+W\coloneqq\{\mathbf u+\mathbf w: \mathbf u\in U,\ \mathbf w\in W\}.
\]
The sum is \term{direct} if $U\cap W=\{\mathbf 0\}$, and it is then written $U\oplus W$. More generally, the sum $U_1+\cdots+U_k\coloneqq\{\mathbf u_1+\cdots+\mathbf u_k: \mathbf u_j\in U_j\}$ of $k$ subspaces is direct, written $U_1\oplus\cdots\oplus U_k$, if $\mathbf u_1+\cdots+\mathbf u_k=\mathbf 0$ with $\mathbf u_j\in U_j$ holds only for $\mathbf u_1=\cdots=\mathbf u_k=\mathbf 0$.
\end{definition}

\begin{theorem}[Sums of subspaces]
\label{thm:direct-sum}
Let $U$ and $W$ be subspaces of $V$.
\begin{enumerate}
\item $U+W$ is a subspace of $V$; it contains $U$ and $W$, and it is contained in every subspace that contains $U$ and $W$.
\item Every vector of $U+W$ can be written as $\mathbf u+\mathbf w$ with $\mathbf u\in U$ and $\mathbf w\in W$ in exactly one way if and only if $U\cap W=\{\mathbf 0\}$.
\item The sum $U_1+\cdots+U_k$ is direct if and only if every vector of it can be written as $\mathbf u_1+\cdots+\mathbf u_k$ with $\mathbf u_j\in U_j$ in exactly one way. For $k=2$ the two definitions of a direct sum agree.
\end{enumerate}
\end{theorem}

\begin{proof}
(i) $\mathbf 0=\mathbf 0+\mathbf 0\in U+W$, and $(\mathbf u+\mathbf w)+(\mathbf u'+\mathbf w')=(\mathbf u+\mathbf u')+(\mathbf w+\mathbf w')$ and $t(\mathbf u+\mathbf w)=t\mathbf u+t\mathbf w$ lie in $U+W$. Taking $\mathbf w=\mathbf 0$ shows $U\subseteq U+W$, and similarly $W\subseteq U+W$. A subspace that contains $U$ and $W$ contains all sums $\mathbf u+\mathbf w$.

(iii) If the decomposition of every vector is unique, it is unique for $\mathbf 0=\mathbf 0+\cdots+\mathbf 0$, which is the definition of a direct sum. Conversely, if $\mathbf u_1+\cdots+\mathbf u_k=\mathbf u_1'+\cdots+\mathbf u_k'$ with $\mathbf u_j,\mathbf u_j'\in U_j$, then $(\mathbf u_1-\mathbf u_1')+\cdots+(\mathbf u_k-\mathbf u_k')=\mathbf 0$ with $\mathbf u_j-\mathbf u_j'\in U_j$, and directness gives $\mathbf u_j=\mathbf u_j'$ for every $j$.

(ii) By (iii), the decomposition is unique if and only if $\mathbf u+\mathbf w=\mathbf 0$ with $\mathbf u\in U$, $\mathbf w\in W$ forces $\mathbf u=\mathbf w=\mathbf 0$. If $U\cap W=\{\mathbf 0\}$ and $\mathbf u+\mathbf w=\mathbf 0$, then $\mathbf u=-\mathbf w$ lies in both $U$ and $W$, so $\mathbf u=\mathbf 0$ and $\mathbf w=\mathbf 0$. Conversely, if $\mathbf v\in U\cap W$, then $\mathbf v+(-\mathbf v)=\mathbf 0$ with $\mathbf v\in U$ and $-\mathbf v\in W$, so uniqueness forces $\mathbf v=\mathbf 0$. This also shows that for $k=2$ the condition of \cref{def:sum} for $U_1\oplus U_2$ is equivalent to $U_1\cap U_2=\{\mathbf 0\}$.
\end{proof}

Direct sums split a space into independent parts. The plane $U=\{\mathbf x\in\mathbb R^3: x_3=0\}$ and the line $W$ spanned by $(1,1,1)$ meet only in $\mathbf 0$, and every vector of $\mathbb R^3$ is the sum
\[
(x_1,x_2,x_3)=(x_1-x_3,\,x_2-x_3,\,0)+x_3(1,1,1)
\]
of a vector of $U$ and one of $W$; so $\mathbb R^3=U\oplus W$. By \cref{prop:symmetric-parts}, every square matrix is in exactly one way the sum of a symmetric and an antisymmetric matrix, so $\mathbb F^{n\times n}$ is the direct sum of the symmetric and the antisymmetric matrices. In the same way, every function $f\colon\mathbb R\to\mathbb R$ is the sum of the even function $(f(x)+f(-x))/2$ and the odd function $(f(x)-f(-x))/2$, and a function that is both even and odd satisfies $f(x)=f(-x)=-f(x)$ and is zero; so the space of all real functions on $\mathbb R$ is the direct sum of the even and the odd functions. For the exponential function the decomposition is $e^x=\cosh x+\sinh x$ (\cref{sec:hyperbolic}).

For more than two subspaces, the condition that the pairwise intersections are $\{\mathbf 0\}$ is not enough for a direct sum. Three different lines $U_1,U_2,U_3$ through the origin of $\mathbb R^2$ meet pairwise only in $\mathbf 0$, but the sum is not direct: with $U_1$, $U_2$, $U_3$ spanned by $(1,0)$, $(0,1)$ and $(1,1)$, the relation $(1,0)+(0,1)-(1,1)=\mathbf 0$ is a decomposition of $\mathbf 0$ with nonzero parts. The correct condition, by the definition, is that $\mathbf 0$ has only the trivial decomposition; equivalently, each $U_j$ meets the sum of the others only in $\mathbf 0$. Direct sums of several subspaces will reappear in Chapter~19, where a matrix that can be diagonalised splits the space into the direct sum of its eigenspaces.

\subsection{Flats and the Solutions of Linear Systems}

A line that does not pass through the origin is not a subspace, but it is a subspace shifted away from the origin, and so is the solution set of every consistent linear system (\cref{thm:solution-structure}).

\begin{definition}[Flat]
\label{def:flat}
For a vector $\mathbf x_0\in V$ and a subspace $U$ of $V$, the set
\[
\mathbf x_0+U\coloneqq\{\mathbf x_0+\mathbf u: \mathbf u\in U\}
\]
is a \term{flat}, or \term{affine subspace}, of $V$, and $U$ is its \term{direction}.
\end{definition}

\begin{proposition}[Flats]
\label{prop:flats}
Let $U$ and $W$ be subspaces of $V$ and $\mathbf x_0,\mathbf x_1\in V$.
\begin{enumerate}
\item $\mathbf x_1\in\mathbf x_0+U$ if and only if $\mathbf x_0+U=\mathbf x_1+U$, and this is the case if and only if $\mathbf x_1-\mathbf x_0\in U$. So every point of a flat can serve as its base point.
\item If $\mathbf x_0+U=\mathbf x_1+W$, then $U=W$: a flat determines its direction, which is the set of all differences of its points.
\item A flat is a subspace if and only if it contains $\mathbf 0$, and then it equals its direction.
\item Two flats with the same direction are either equal or disjoint.
\end{enumerate}
\end{proposition}

\begin{proof}
(i) If $\mathbf x_1=\mathbf x_0+\mathbf u$ with $\mathbf u\in U$, then $\mathbf x_1+\mathbf u'=\mathbf x_0+(\mathbf u+\mathbf u')$ and $\mathbf x_0+\mathbf u'=\mathbf x_1+(\mathbf u'-\mathbf u)$ for every $\mathbf u'\in U$, so the two flats coincide. Conversely, if $\mathbf x_0+U=\mathbf x_1+U$, then $\mathbf x_1=\mathbf x_1+\mathbf 0$ lies in $\mathbf x_0+U$, so $\mathbf x_1-\mathbf x_0\in U$. (ii) The differences of the points of $\mathbf x_0+U$ are the vectors $(\mathbf x_0+\mathbf u)-(\mathbf x_0+\mathbf u')=\mathbf u-\mathbf u'$, which are exactly the vectors of $U$; the differences of the points of $\mathbf x_1+W$ are the vectors of $W$, and the two sets of differences are the same. (iii) If $\mathbf 0\in\mathbf x_0+U$, then $\mathbf x_0+U=\mathbf 0+U=U$ by (i); conversely, a subspace contains $\mathbf 0$. (iv) If $\mathbf x_0+U$ and $\mathbf x_1+U$ have a common point $\mathbf y$, then both equal $\mathbf y+U$ by (i).
\end{proof}

The lines and planes of \cref{sec:vectors-in-rn} are the flats $\mathbf p+\operatorname{span}\{\mathbf v\}$ and $\mathbf p+\operatorname{span}\{\mathbf u,\mathbf v\}$, and \cref{thm:solution-structure} now reads: \emph{the solution set of a consistent system $A\mathbf x=\mathbf b$ is the flat $\mathbf x_{\mathrm p}+\mathcal N(A)$, whose direction is the null space of $A$}. By (i), every solution can serve as the particular solution, and by (iv), the solution sets of the systems $A\mathbf x=\mathbf b$ for the different right-hand sides $\mathbf b\in\mathcal C(A)$ are parallel flats that do not meet, together filling $\mathbb F^n$. A flat that is not a subspace is never closed under the operations: the sum of two points of $\mathbf x_0+U$ with $\mathbf x_0\notin U$ lies in the different flat $2\mathbf x_0+U$. \Cref{fig:subspaces} shows subspaces and flats in $\mathbb R^3$.

\begin{figure}
\newcommand\sdthreeaxes{%
\draw[line width=0.4pt, -{Latex[length=1.5mm]}] (0,0,0) -- (2.3,0,0) node[below left=-1pt] {$x_1$};
\draw[line width=0.4pt, -{Latex[length=1.5mm]}] (0,0,0) -- (0,2.6,0) node[right] {$x_2$};
\draw[line width=0.4pt, -{Latex[length=1.5mm]}] (0,0,0) -- (0,0,2.4) node[above] {$x_3$};}
\sdpanel{A}{\begin{tikzpicture}[x={(-0.42cm,-0.36cm)}, y={(1cm,0cm)}, z={(0cm,0.9cm)}, font=\scriptsize]
\fill[sdfill] (-1.6,-1.4,0) -- (1.9,-1.4,0) -- (1.9,2.0,0) -- (-1.6,2.0,0) -- cycle;
\draw[black!50, line width=0.4pt] (-1.6,-1.4,0) -- (1.9,-1.4,0) -- (1.9,2.0,0) -- (-1.6,2.0,0) -- cycle;
\draw[black!60, dashed, line width=0.5pt] (-1.6,-1.4,1.6) -- (1.9,-1.4,1.6) -- (1.9,2.0,1.6) -- (-1.6,2.0,1.6) -- cycle;
\sdthreeaxes
\draw[line width=0.9pt] (-0.9,-0.9,-0.9) -- (1.2,1.2,1.2);
\node[right] at (1.2,1.2,1.2) {$W$};
\node at (1.4,1.6,0) {$U$};
\node[right] at (-1.6,2.0,1.6) {$\mathbf x_0+U$};
\fill (0,0,1.6) circle (1.6pt);
\node[left] at (0,0,1.6) {$\mathbf x_0$};
\fill (0,0,0) circle (1.4pt);
\end{tikzpicture}}\hfill
\sdpanel{B}{\begin{tikzpicture}[x={(-0.42cm,-0.36cm)}, y={(1cm,0cm)}, z={(0cm,0.9cm)}, font=\scriptsize]
\fill[sdfill] (-1.6,-1.4,0) -- (1.9,-1.4,0) -- (1.9,2.0,0) -- (-1.6,2.0,0) -- cycle;
\draw[black!50, line width=0.4pt] (-1.6,-1.4,0) -- (1.9,-1.4,0) -- (1.9,2.0,0) -- (-1.6,2.0,0) -- cycle;
\fill[black!5, opacity=0.8] (0,-1.4,-1.2) -- (0,2.0,-1.2) -- (0,2.0,1.8) -- (0,-1.4,1.8) -- cycle;
\draw[black!50, line width=0.4pt] (0,-1.4,-1.2) -- (0,2.0,-1.2) -- (0,2.0,1.8) -- (0,-1.4,1.8) -- cycle;
\sdthreeaxes
\draw[line width=1.3pt] (0,-1.4,0) -- (0,2.0,0);
\node at (1.4,1.6,0) {$U_1$};
\node at (0,1.6,1.4) {$U_2$};
\node[above left] at (0,2.0,0) {$U_1\cap U_2$};
\end{tikzpicture}}
\caption{Subspaces and flats of $\mathbb R^3$. (A) The plane $U=\{\mathbf x: x_3=0\}$ (shaded) and the line $W$ spanned by $(1,1,1)$ are subspaces, with $U\cap W=\{\mathbf 0\}$ and $U\oplus W=\mathbb R^3$. The parallel plane $\mathbf x_0+U$ through $\mathbf x_0=(0,0,1)$ (dashed) is a flat with direction $U$, the solution set of the equation $x_3=1$, and it is not a subspace. (B) Two different planes $U_1=\{\mathbf x: x_3=0\}$ and $U_2=\{\mathbf x: x_1=0\}$ through the origin meet in the line $U_1\cap U_2$, the $x_2$-axis, which is again a subspace.}
\label{fig:subspaces}
\end{figure}

\section{Span and Linear Independence}
\label{sec:independence}

\subsection{Span}

The span of finitely many vectors of $\mathbb R^n$ was defined in \cref{def:span}, and the definition carries over to every vector space. For infinite sets we allow only finite combinations, since an infinite sum of vectors has no meaning in a vector space without a notion of limit.

\begin{definition}[Span in a vector space]
\label{def:span-general}
Let $V$ be a vector space over $\mathbb F$. A \term{linear combination} of $\mathbf v_1,\ldots,\mathbf v_k\in V$ is a vector $t_1\mathbf v_1+\cdots+t_k\mathbf v_k$ with $t_1,\ldots,t_k\in\mathbb F$. The \term{span} of a subset $S\subseteq V$, written $\operatorname{span}S$, is the set of all linear combinations of finitely many elements of $S$, and $\operatorname{span}\emptyset\coloneqq\{\mathbf 0\}$. If $\operatorname{span}S=V$, then $S$ \term{spans} $V$, or \term{generates} $V$, and $V$ is \term{finitely generated} if it is spanned by a finite set. For a list of vectors we write $\operatorname{span}\{\mathbf v_1,\ldots,\mathbf v_k\}$ as before.
\end{definition}

\begin{proposition}[The span is the smallest subspace]
\label{prop:span-smallest}
For every subset $S$ of a vector space $V$, $\operatorname{span}S$ is a subspace of $V$ that contains $S$, and it is contained in every subspace of $V$ that contains $S$. Consequently $\operatorname{span}S$ is the intersection of all subspaces of $V$ that contain $S$.
\end{proposition}

\begin{proof}
If $S$ is empty, the statements are clear, since every subspace contains $\{\mathbf 0\}$. Otherwise $\operatorname{span}S$ contains every $\mathbf s\in S$ as the combination $1\mathbf s$, and $\mathbf 0=0\mathbf s$. A sum of two finite combinations of elements of $S$ is again one, with the coefficients of an element that occurs in both added, and a multiple $t\sum_js_j\mathbf v_j=\sum_j(ts_j)\mathbf v_j$ is one as well. By the subspace test, $\operatorname{span}S$ is a subspace. A subspace $U$ that contains $S$ contains all multiples $t\mathbf s$ with $\mathbf s\in S$ and, by induction on the number of terms, all finite sums of such multiples; so $\operatorname{span}S\subseteq U$. The intersection of all subspaces containing $S$ is one of them, by \cref{prop:intersection}, and it lies inside each of them, in particular inside $\operatorname{span}S$; conversely it contains $\operatorname{span}S$ by what was just shown.
\end{proof}

This is the statement that \cref{prop:span} proved in $\mathbb R^n$, with the word that \cref{ch:vectors-matrices} anticipated: the span of $S$ is the smallest subspace containing $S$. The two descriptions in the proposition are typical of two ways of producing an object: from inside, by building all combinations, and from outside, by intersecting everything that contains the given set. The examples of \cref{sec:subspaces} have natural spanning sets. The monomials $1,x,\ldots,x^n$ span $\mathcal P_n$, by the definition of a polynomial, and the infinite set $\{1,x,x^2,\ldots\}$ spans $\mathcal P$. The $mn$ \term{matrix units} $E_{ij}\coloneqq\mathbf e_i\mathbf e_j^{\mathsf T}$, with the entry $1$ in position $(i,j)$ and zeros elsewhere, span $\mathbb F^{m\times n}$, since $A=\sum_{i,j}a_{ij}E_{ij}$. The functions $\cos\omega x$ and $\sin\omega x$ span the solutions of $f''=-\omega^2f$, by \cref{sec:vector-spaces}. The columns of $A$ span $\mathcal C(A)$, by definition.

Adding a vector to a spanning list enlarges the span only if the vector is new, a fact that is used again and again.

\begin{proposition}[Enlarging a span]
\label{prop:span-enlarge}
Let $\mathbf v_1,\ldots,\mathbf v_k,\mathbf w\in V$. Then
\[
\operatorname{span}\{\mathbf v_1,\ldots,\mathbf v_k,\mathbf w\}=\operatorname{span}\{\mathbf v_1,\ldots,\mathbf v_k\}
\]
if and only if $\mathbf w\in\operatorname{span}\{\mathbf v_1,\ldots,\mathbf v_k\}$.
\end{proposition}

\begin{proof}
If the spans are equal, $\mathbf w$ lies in the second one. Conversely, if $\mathbf w$ lies in the subspace $\operatorname{span}\{\mathbf v_1,\ldots,\mathbf v_k\}$, this subspace contains all of $\mathbf v_1,\ldots,\mathbf v_k,\mathbf w$, hence their span, by \cref{prop:span-smallest}; the reverse inclusion is clear.
\end{proof}

\subsection{Linear Independence}

In the plane, two vectors that are not multiples of each other span the whole plane, while two vectors on a common line span only the line; the second vector then adds nothing. Linear independence makes the idea of ``adding nothing new'' precise.

\begin{definition}[Linear independence]
\label{def:independence}
Vectors $\mathbf v_1,\ldots,\mathbf v_k$ of a vector space $V$ are \term{linearly independent} if
\[
t_1\mathbf v_1+t_2\mathbf v_2+\cdots+t_k\mathbf v_k=\mathbf 0\qquad\text{only for}\qquad t_1=t_2=\cdots=t_k=0,
\]
and \term{linearly dependent} otherwise; a combination equal to $\mathbf 0$ with coefficients not all zero is a \term{linear relation} among the vectors. The empty list is independent. An infinite subset $S\subseteq V$ is linearly independent if every finite list of distinct elements of $S$ is independent.
\end{definition}

Independence is a property of the list as a whole, not of its members. A list that contains the zero vector is dependent, because $1\cdot\mathbf 0=\mathbf 0$ is a relation, and so is a list in which a vector occurs twice, because $\mathbf v-\mathbf v=\mathbf 0$. A single vector is independent exactly when it is nonzero, by \cref{prop:vector-space-rules}\,(vi), and two vectors are dependent exactly when one of them is a multiple of the other. Every sublist of an independent list is independent, since a relation among some of the vectors is a relation among all of them, with zero coefficients for the others. The following theorem collects the equivalent forms of the definition.

\begin{theorem}[Forms of independence]
\label{thm:independence}
For vectors $\mathbf v_1,\ldots,\mathbf v_k$ of $V$ with $k\ge1$ the following are equivalent:
\begin{enumerate}
\item $\mathbf v_1,\ldots,\mathbf v_k$ are linearly independent;
\item none of the vectors is a linear combination of the others;
\item every vector of $\operatorname{span}\{\mathbf v_1,\ldots,\mathbf v_k\}$ can be written as $t_1\mathbf v_1+\cdots+t_k\mathbf v_k$ in exactly one way;
\item $\mathbf v_1\ne\mathbf 0$, and $\mathbf v_j\notin\operatorname{span}\{\mathbf v_1,\ldots,\mathbf v_{j-1}\}$ for $j=2,\ldots,k$.
\end{enumerate}
\end{theorem}

\begin{proof}
(i)$\Rightarrow$(iii): If $\sum_jt_j\mathbf v_j=\sum_js_j\mathbf v_j$, then $\sum_j(t_j-s_j)\mathbf v_j=\mathbf 0$, and independence gives $t_j=s_j$ for every $j$. (iii)$\Rightarrow$(i): The zero vector has the representation with all coefficients $0$, and by (iii) no other.

(i)$\Rightarrow$(ii): If $\mathbf v_j=\sum_{i\ne j}c_i\mathbf v_i$, then $\sum_{i\ne j}c_i\mathbf v_i-\mathbf v_j=\mathbf 0$ is a relation, with the coefficient $-1$ for $\mathbf v_j$. (ii)$\Rightarrow$(i): If $\sum_it_i\mathbf v_i=\mathbf 0$ with some $t_j\ne0$, then $\mathbf v_j=-\sum_{i\ne j}(t_i/t_j)\mathbf v_i$ is a combination of the others.

(i)$\Rightarrow$(iv): (iv) says that no vector is a combination of those before it, which is a special case of (ii), and that $\mathbf v_1\ne\mathbf 0$, which holds because a list containing $\mathbf 0$ is dependent. (iv)$\Rightarrow$(i): Suppose that $\sum_it_i\mathbf v_i=\mathbf 0$ with coefficients not all zero, and let $j$ be the largest index with $t_j\ne0$. If $j=1$, then $t_1\mathbf v_1=\mathbf 0$ and $\mathbf v_1=\mathbf 0$. If $j\ge2$, then $\mathbf v_j=-\sum_{i<j}(t_i/t_j)\mathbf v_i$ lies in the span of $\mathbf v_1,\ldots,\mathbf v_{j-1}$. Either way (iv) fails.
\end{proof}

Form (iii) is the one that makes independence useful: the coefficients of a vector with respect to independent vectors are well defined, and they can serve as coordinates (\cref{sec:coordinates}). Form (iv) is the one that is used to build independent lists one vector at a time, and it gives at once the following lemma.

\begin{lemma}[Extending an independent list]
\label{lem:extend-independent}
If $\mathbf v_1,\ldots,\mathbf v_k$ are linearly independent and $\mathbf w\notin\operatorname{span}\{\mathbf v_1,\ldots,\mathbf v_k\}$, then $\mathbf v_1,\ldots,\mathbf v_k,\mathbf w$ are linearly independent.
\end{lemma}

\begin{proof}
The list $\mathbf v_1,\ldots,\mathbf v_k$ satisfies condition (iv) of \cref{thm:independence}, and the new vector $\mathbf w$ does not lie in the span of its predecessors; so the longer list satisfies (iv) as well. (For $k=0$, $\mathbf w\notin\{\mathbf 0\}$ means $\mathbf w\ne\mathbf 0$.)
\end{proof}

\Cref{fig:independence} shows the possibilities in the plane. Two vectors that are not multiples of each other are independent and span $\mathbb R^2$; two vectors on a line are dependent; and three vectors in the plane are always dependent, as the next proposition shows.

\begin{figure}
\newcommand\sdplaneaxes{%
\draw[line width=0.4pt, -{Latex[length=1.5mm]}] (-1.1,0) -- (3.3,0) node[below] {$x_1$};
\draw[line width=0.4pt, -{Latex[length=1.5mm]}] (0,-1.3) -- (0,3.3) node[left] {$x_2$};}
\sdpanel{A}{\begin{tikzpicture}[scale=0.82, font=\scriptsize]
\fill[sdfill] (0,0) -- (2,1) -- (3,3) -- (1,2) -- cycle;
\sdplaneaxes
\draw[line width=0.8pt, -{Latex[length=2mm]}] (0,0) -- (2,1) node[below right=-1pt] {$\mathbf u$};
\draw[line width=0.8pt, -{Latex[length=2mm]}] (0,0) -- (1,2) node[above left=-1pt] {$\mathbf v$};
\end{tikzpicture}}\hfill
\sdpanel{B}{\begin{tikzpicture}[scale=0.82, font=\scriptsize]
\draw[black!45, dashed, line width=0.5pt] (-0.6,-1.2) -- (1.6,3.2);
\sdplaneaxes
\draw[line width=1.2pt, -{Latex[length=2mm]}] (0,0) -- (1,2) node[right] {$\mathbf w=2\mathbf u$};
\draw[line width=0.8pt, white, -{Latex[length=2mm]}] (0,0) -- (0.5,1);
\draw[line width=0.8pt, -{Latex[length=2mm]}] (0,0) -- (0.5,1) node[left] {$\mathbf u$};
\end{tikzpicture}}\hfill
\sdpanel{C}{\begin{tikzpicture}[scale=0.82, font=\scriptsize]
\sdplaneaxes
\draw[black!60, dashed, line width=0.5pt, -{Latex[length=1.6mm]}] (2,1) -- (1,-1);
\draw[line width=0.8pt, -{Latex[length=2mm]}] (0,0) -- (2,1) node[above right=-1pt] {$\mathbf u$};
\draw[line width=0.8pt, -{Latex[length=2mm]}] (0,0) -- (1,2) node[above left=-1pt] {$\mathbf v$};
\draw[line width=1.2pt, -{Latex[length=2mm]}] (0,0) -- (1,-1) node[below right=-1pt] {$\mathbf w=\mathbf u-\mathbf v$};
\end{tikzpicture}}
\caption{Independence in the plane. (A) The vectors $\mathbf u=(2,1)$ and $\mathbf v=(1,2)$ are independent; the parallelogram they span (shaded) does not collapse to a segment, and their combinations fill the plane. (B) The vectors $\mathbf u=(1/2,1)$ and $\mathbf w=2\mathbf u=(1,2)$ lie on one line and are dependent; their span is the line (dashed). (C) Three vectors in the plane are always dependent: $\mathbf w=(1,-1)$ is the combination $\mathbf u-\mathbf v$ (dashed: $\mathbf u$ followed by $-\mathbf v$), so $\mathbf u-\mathbf v-\mathbf w=\mathbf 0$ is a linear relation.}
\label{fig:independence}
\end{figure}

\subsection{Independence in \texorpdfstring{$\mathbb F^m$}{F\^{}m}}

For vectors with numerical components, independence is decided by elimination. If $\mathbf a_1,\ldots,\mathbf a_n\in\mathbb F^m$ are the columns of $A$, a combination $\sum_jx_j\mathbf a_j$ is the product $A\mathbf x$, and a linear relation among the columns is a nonzero solution of $A\mathbf x=\mathbf 0$.

\begin{proposition}[Independence of columns]
\label{prop:independence-columns}
Let $A\in\mathbb F^{m\times n}$ have the columns $\mathbf a_1,\ldots,\mathbf a_n$.
\begin{enumerate}
\item The columns are linearly independent if and only if $\mathcal N(A)=\{\mathbf 0\}$, that is, if and only if $\operatorname{rank}A=n$.
\item The columns span $\mathbb F^m$ if and only if $\operatorname{rank}A=m$.
\item If $n>m$, the columns are linearly dependent: more than $m$ vectors in $\mathbb F^m$ are always dependent.
\item If $n<m$, the columns do not span $\mathbb F^m$.
\end{enumerate}
\end{proposition}

\begin{proof}
(i) The relations $\sum_jx_j\mathbf a_j=\mathbf 0$ are the solutions of $A\mathbf x=\mathbf 0$, and by \cref{thm:homogeneous-solutions} there is a nonzero one if and only if there is a free variable, that is, $\operatorname{rank}A<n$. (ii) The columns span $\mathbb F^m$ if and only if $A\mathbf x=\mathbf b$ is consistent for every $\mathbf b$, which by \cref{thm:rank-solutions}\,(iii) means $\operatorname{rank}A=m$. (iii) is \cref{thm:more-unknowns}, and (iv) follows from (ii), since $\operatorname{rank}A\le n<m$ by \cref{prop:rank-basic}\,(i).
\end{proof}

The elimination also produces the relations. For the vectors $(1,2,3)$, $(4,5,6)$ and $(7,8,9)$, the matrix with these columns reduces to
\[
A=\begin{bmatrix}1&4&7\\2&5&8\\3&6&9\end{bmatrix}\;\longrightarrow\;R=\begin{bmatrix}1&0&-1\\0&1&2\\0&0&0\end{bmatrix},
\]
with the free column $3$ and the special solution $(1,-2,1)$. So the vectors are dependent, and the relation is $(1,2,3)-2(4,5,6)+(7,8,9)=\mathbf 0$, as one checks at once. The third vector lies in the plane spanned by the first two.

The proposition settles the question of \cref{sec:vectors-in-rn} about spans in $\mathbb R^3$. The vectors $\mathbf v_1=(1,0,1)$ and $\mathbf v_2=(0,1,1)$ are independent and span the plane $x_3=x_1+x_2$. For a third vector $\mathbf w$, \cref{lem:extend-independent} and \cref{thm:independence}\,(iv) show that $\mathbf v_1,\mathbf v_2,\mathbf w$ are independent if and only if $\mathbf w$ does not lie in the plane. If they are independent, the $3\times3$ matrix with these columns has rank $3$ by (i), and then its columns span $\mathbb R^3$ by (ii). If they are dependent, $\mathbf w$ lies in the plane, and the span stays the plane. So a third vector enlarges the span to all of $\mathbb R^3$ exactly when it is independent of the first two, and otherwise it enlarges nothing. Panel~C of \cref{fig:independence} is the case of $\mathbb R^2$ in part (iii): three vectors in the plane are always dependent.

\subsection{Independence of Functions}

In spaces of functions, a relation $t_1f_1+\cdots+t_kf_k=0$ means that the combination vanishes at \emph{every} point, and to show that functions are independent one exploits this abundance of conditions.

The monomials $1,x,x^2,\ldots,x^n$ are independent in the space of functions on $\mathbb R$ or on any interval with more than one point, since a polynomial $c_0+c_1x+\cdots+c_nx^n$ that vanishes at infinitely many points has all coefficients zero (\cref{cor:factor}). Consequently the infinite set $\{1,x,x^2,\ldots\}$ is independent as well, since each of its finite subsets is contained in some list $1,x,\ldots,x^n$. The functions $\cos x$ and $\sin x$ are independent: if $a\cos x+b\sin x=0$ for all $x$, then $x=0$ gives $a=0$ and $x=\pi/2$ gives $b=0$. The three functions $\sin x$, $\cos x$ and $\sin(x+1)$ are dependent, since $\sin(x+1)=\cos1\,\sin x+\sin1\,\cos x$ by the addition formula (\cref{thm:addition}).

\begin{proposition}[Exponentials are independent]
\label{prop:exponentials-independent}
For pairwise different real numbers $\lambda_1,\ldots,\lambda_k$, the functions $e^{\lambda_1x},\ldots,e^{\lambda_kx}$ are linearly independent in the space of real functions on any interval $I$ with more than one point.
\end{proposition}

\begin{proof}
By induction on $k$. For $k=1$, $e^{\lambda_1x}\ne0$. Let $k\ge2$, and suppose that $\sum_{j=1}^kc_je^{\lambda_jx}=0$ for all $x\in I$. Differentiating gives $\sum_jc_j\lambda_je^{\lambda_jx}=0$, and subtracting $\lambda_k$ times the original relation eliminates the last term:
\[
\sum_{j=1}^{k-1}c_j(\lambda_j-\lambda_k)e^{\lambda_jx}=0\qquad(x\in I).
\]
By the induction hypothesis, $c_j(\lambda_j-\lambda_k)=0$ for $j<k$, and since $\lambda_j\ne\lambda_k$, $c_j=0$. The original relation reduces to $c_ke^{\lambda_kx}=0$, so $c_k=0$.
\end{proof}

The proof is elimination in disguise: differentiation acts on each exponential by multiplication with its own number $\lambda_j$, and subtracting a multiple of the relation removes one exponential, as a row operation removes one unknown. A second proof uses a general test, which deserves to be stated separately because it reduces the independence of functions to the independence of vectors of numbers.

\begin{proposition}[Independence from derivatives at a point]
\label{prop:derivative-test}
Let $f_1,\ldots,f_k$ be $k-1$ times differentiable on an interval $I$, let $x_0\in I$, and let
\[
\mathbf w_j\coloneqq\bigl(f_j(x_0),\,f_j'(x_0),\,\ldots,\,f_j^{(k-1)}(x_0)\bigr)\in\mathbb R^k\qquad(j=1,\ldots,k).
\]
If $\mathbf w_1,\ldots,\mathbf w_k$ are linearly independent, then so are $f_1,\ldots,f_k$.
\end{proposition}

\begin{proof}
If $\sum_jc_jf_j=0$ on $I$, then also $\sum_jc_jf_j^{(i)}=0$ on $I$ for $i=1,\ldots,k-1$, by differentiating $i$ times. At $x_0$ these $k$ equations say that $\sum_jc_j\mathbf w_j=\mathbf 0$, so all $c_j$ vanish.
\end{proof}

For the exponentials and $x_0=0$, the vector $\mathbf w_j$ is $(1,\lambda_j,\lambda_j^2,\ldots,\lambda_j^{k-1})$, and the matrix with the columns $\mathbf w_1,\ldots,\mathbf w_k$ is the transpose of the Vandermonde matrix of the points $\lambda_1,\ldots,\lambda_k$ (\cref{sec:special-systems}). It is invertible by \cref{thm:interpolation}, since the points are different, and so is its transpose (\cref{prop:inverse-rules}\,(iii)); so its columns are independent, by \cref{thm:invertible}, and the proposition gives a second proof of \cref{prop:exponentials-independent}. The converse of \cref{prop:derivative-test} is false. The functions $f_1(x)=x^2$ and $f_2(x)=x|x|$ are differentiable on $\mathbb R$, with $f_1'(x)=2x$ and $f_2'(x)=2|x|$, and for every $x_0$ the vectors $(f_1(x_0),f_1'(x_0))$ and $(f_2(x_0),f_2'(x_0))$ are dependent: they are equal for $x_0\ge0$ and opposite for $x_0<0$. Yet the functions are independent on $\mathbb R$, since $c_1x^2+c_2x|x|=0$ gives $c_1+c_2=0$ at $x=1$ and $c_1-c_2=0$ at $x=-1$. They are proportional on each half-line, but with different factors.

The test is usually phrased with the matrix whose columns are $\mathbf w_1,\ldots,\mathbf w_k$, regarded as a function of $x_0$; its determinant is the \term{Wronskian} of $f_1,\ldots,f_k$, and the test says that the functions are independent if the Wronskian is nonzero at a single point. Determinants are defined in Chapter~17. For solutions of a linear differential equation of order $k$, the example above cannot occur, and the Wronskian vanishes either everywhere or nowhere; Chapter~24 proves this and uses it to show that the solutions of such an equation form a space of dimension $k$.

\section{Bases and Dimension}
\label{sec:basis-dimension}

\subsection{Bases}

A spanning list produces every vector of the space, and an independent list produces each of them in at most one way. A list that does both is the most economical description of a space.

\begin{definition}[Basis]
\label{def:basis}
A \term{basis} of a vector space $V$ is a list $(\mathbf b_1,\ldots,\mathbf b_n)$ of vectors of $V$ that is linearly independent and spans $V$. More generally, a subset of $V$ that is linearly independent and spans $V$ is called a basis as well. The empty list is a basis of $\{\mathbf 0\}$.
\end{definition}

\begin{theorem}[Bases and unique representation]
\label{thm:basis-unique}
A list $(\mathbf b_1,\ldots,\mathbf b_n)$ is a basis of $V$ if and only if every $\mathbf v\in V$ can be written as $\mathbf v=t_1\mathbf b_1+\cdots+t_n\mathbf b_n$ with scalars $t_1,\ldots,t_n$ in exactly one way.
\end{theorem}

\begin{proof}
That every vector can be written in this way at all means that the list spans $V$, and that it can be written in only one way means, by \cref{thm:independence}\,(iii), that the list is independent.
\end{proof}

The \term{standard bases} of the examples are the following. In $\mathbb F^n$, the standard basis vectors $\mathbf e_1,\ldots,\mathbf e_n$ form a basis, and the coefficients of $\mathbf x$ are its components, $\mathbf x=\sum_ix_i\mathbf e_i$. In $\mathbb F^{m\times n}$, the matrix units $E_{ij}$ form a basis, and the coefficients of $A$ are its entries. In $\mathcal P_n$, the monomials $1,x,\ldots,x^n$ form a basis, and the coefficients of a polynomial are its coefficients in the usual sense; they span by definition and are independent by \cref{sec:independence}. The infinite set $\{1,x,x^2,\ldots\}$ is a basis of $\mathcal P$. In the space of solutions of $f''=-\omega^2f$, the functions $\cos\omega x$ and $\sin\omega x$ form a basis: they span by \cref{lem:harmonic-unique}, as explained in \cref{sec:vector-spaces}, and they are independent by evaluating a relation at $x=0$ and $x=\pi/(2\omega)$. In the real vector space $\mathbb C^n$, the $2n$ vectors $\mathbf e_1,\ldots,\mathbf e_n,i\mathbf e_1,\ldots,i\mathbf e_n$ form a basis, since every vector is uniquely $\sum_ja_j\mathbf e_j+\sum_jb_j(i\mathbf e_j)$ with the real and imaginary parts $a_j$ and $b_j$ of its components; as a complex vector space, $\mathbb C^n$ has the basis $\mathbf e_1,\ldots,\mathbf e_n$.

A vector space has many bases. Any two vectors of $\mathbb R^2$ that are not multiples of each other form a basis, by \cref{prop:independence-columns}, since the $2\times2$ matrix with these columns then has two independent columns and hence rank $2$. What all bases of a space have in common is the number of their elements. This is the main theorem of the chapter, and it rests on the following lemma.

\subsection{The Exchange Lemma and the Dimension}

\begin{lemma}[Exchange lemma]
\label{lem:steinitz}
Let $\mathbf w_1,\ldots,\mathbf w_m$ span $V$, and let $\mathbf v_1,\ldots,\mathbf v_k\in V$ be linearly independent. Then:
\begin{enumerate}
\item $k\le m$;
\item after renumbering $\mathbf w_1,\ldots,\mathbf w_m$, the list $\mathbf v_1,\ldots,\mathbf v_k$, $\mathbf w_{k+1},\ldots,\mathbf w_m$ spans $V$.
\end{enumerate}
\end{lemma}

\begin{proof}
(i) Since the $\mathbf w_i$ span $V$, each $\mathbf v_j$ is a combination of them, $\mathbf v_j=\sum_{i=1}^ma_{ij}\mathbf w_i$, with a matrix $A=(a_{ij})\in\mathbb F^{m\times k}$ of coefficients. Suppose that $k>m$. Then the homogeneous system $A\mathbf x=\mathbf 0$ has more unknowns than equations, and by \cref{thm:more-unknowns} it has a solution $\mathbf x\ne\mathbf 0$. For this $\mathbf x$,
\[
\sum_{j=1}^kx_j\mathbf v_j=\sum_{j=1}^kx_j\sum_{i=1}^ma_{ij}\mathbf w_i=\sum_{i=1}^m\Bigl(\sum_{j=1}^ka_{ij}x_j\Bigr)\mathbf w_i=\sum_{i=1}^m(A\mathbf x)_i\,\mathbf w_i=\mathbf 0,
\]
a linear relation among the $\mathbf v_j$ with coefficients not all zero, contrary to their independence. So $k\le m$.

(ii) We show by induction on $j=0,1,\ldots,k$ that, after renumbering, the list $\mathbf v_1,\ldots,\mathbf v_j$, $\mathbf w_{j+1},\ldots,\mathbf w_m$ spans $V$. For $j=0$ this is the hypothesis. Suppose it holds for some $j<k$. Then
\[
\mathbf v_{j+1}=\sum_{i=1}^ja_i\mathbf v_i+\sum_{i=j+1}^mb_i\mathbf w_i
\]
for some scalars. Not all $b_i$ are zero, since otherwise $\mathbf v_{j+1}$ would lie in the span of $\mathbf v_1,\ldots,\mathbf v_j$, contradicting independence (\cref{thm:independence}\,(iv)); in particular $j+1\le m$. Renumber $\mathbf w_{j+1},\ldots,\mathbf w_m$ so that $b_{j+1}\ne0$. Then
\[
\mathbf w_{j+1}=\frac1{b_{j+1}}\Bigl(\mathbf v_{j+1}-\sum_{i=1}^ja_i\mathbf v_i-\sum_{i=j+2}^mb_i\mathbf w_i\Bigr)
\]
lies in the span of $\mathbf v_1,\ldots,\mathbf v_{j+1},\mathbf w_{j+2},\ldots,\mathbf w_m$. By \cref{prop:span-enlarge}, adding $\mathbf w_{j+1}$ to this list does not change its span; but the enlarged list contains the spanning list of the induction hypothesis, so it spans $V$. Hence $\mathbf v_1,\ldots,\mathbf v_{j+1},\mathbf w_{j+2},\ldots,\mathbf w_m$ spans $V$.
\end{proof}

Part (ii) explains the name: the independent vectors can be exchanged, one at a time, for vectors of the spanning list, without losing the property of spanning. Part (i) is the statement that matters most, and its proof is where the theory of linear systems enters: \emph{the length of an independent list never exceeds the length of a spanning list}. In $\mathbb F^m$, with the standard basis as spanning list, it is \cref{prop:independence-columns}\,(iii) once more.

\begin{theorem}[Dimension]
\label{thm:dimension}
Let $V$ be finitely generated. Then $V$ has a basis, every spanning list of $V$ contains a basis, and all bases of $V$ have the same number of elements.
\end{theorem}

\begin{proof}
Let $\mathbf w_1,\ldots,\mathbf w_m$ span $V$. Go through the list from the first vector to the last, and keep $\mathbf w_j$ if it does not lie in the span of the vectors kept before it; discard it otherwise. The kept vectors are independent, by \cref{thm:independence}\,(iv). Each discarded vector lies in the span of the vectors kept before it, so by \cref{prop:span-enlarge}, applied once for each discarded vector, the kept vectors span the same space as all the $\mathbf w_j$, namely $V$. So the kept vectors form a basis contained in the spanning list. If $(\mathbf b_1,\ldots,\mathbf b_n)$ and $(\mathbf b_1',\ldots,\mathbf b_{n'}')$ are two bases, then the second is independent and the first spans, so $n'\le n$ by \cref{lem:steinitz}; exchanging the roles gives $n\le n'$.
\end{proof}

\begin{definition}[Dimension]
\label{def:dimension}
The \term{dimension} $\dim V$ of a finitely generated vector space $V$ is the number of elements of any of its bases, and $V$ is then called \term{finite-dimensional}. In particular $\dim\{\mathbf 0\}=0$. A vector space that is not finitely generated is \term{infinite-dimensional}.
\end{definition}

So $\dim\mathbb F^n=n$, $\dim\mathbb F^{m\times n}=mn$ and $\dim\mathcal P_n=n+1$, and the solutions of $f''=-\omega^2f$ form a space of dimension $2$. The real vector space $\mathbb C^n$ has dimension $2n$ and the complex vector space $\mathbb C^n$ dimension $n$: the dimension depends on the field of scalars, and when the field is not clear from the context, one writes $\dim_{\mathbb R}$ or $\dim_{\mathbb C}$. In the space $\mathbb R_{>0}$ of \cref{sec:vector-spaces}, every vector other than the zero vector $1$ forms a basis by itself, for instance the number $e$, since $x=e^{\ln x}=(\ln x)\odot e$ for every $x>0$, in exactly one way; so this space has dimension $1$.

\begin{theorem}[Extending and reducing to a basis]
\label{thm:basis-extension}
Let $V$ have dimension $n$.
\begin{enumerate}
\item Every linearly independent list in $V$ has at most $n$ elements, and every spanning list has at least $n$ elements.
\item Every linearly independent list in $V$ can be extended to a basis of $V$, using vectors of any given basis.
\item A linearly independent list of $n$ vectors is a basis of $V$, and so is a spanning list of $n$ vectors.
\end{enumerate}
\end{theorem}

\begin{proof}
(i) A basis has $n$ elements and both spans and is independent; \cref{lem:steinitz}\,(i), applied with the basis as the spanning list or as the independent list, gives both statements.

(ii) Let $\mathbf v_1,\ldots,\mathbf v_k$ be independent and $\mathbf b_1,\ldots,\mathbf b_n$ a basis. The list $\mathbf v_1,\ldots,\mathbf v_k$, $\mathbf b_1,\ldots,\mathbf b_n$ spans $V$. Apply the procedure of the proof of \cref{thm:dimension} to it. None of the $\mathbf v_j$ is discarded, since none lies in the span of its predecessors, by \cref{thm:independence}\,(iv). The result is a basis that begins with $\mathbf v_1,\ldots,\mathbf v_k$ and continues with some of the $\mathbf b_i$.

(iii) An independent list of $n$ vectors extends to a basis by (ii), which has $n$ elements, so nothing is added and the list is already a basis. A spanning list of $n$ vectors contains a basis by \cref{thm:dimension}, which has $n$ elements, so nothing is removed.
\end{proof}

Part (iii) halves the work of checking a basis when the dimension is known: of the two conditions, independence and spanning, one implies the other. For instance, $1$, $1+x$ and $(1+x)^2$ are three polynomials in the space $\mathcal P_2$ of dimension $3$, and they are independent by \cref{thm:independence}\,(iv), since each has a higher degree than the ones before it and so does not lie in their span; hence they form a basis of $\mathcal P_2$, and every polynomial of degree at most $2$ can be written as $a+b(1+x)+c(1+x)^2$, which is its Taylor expansion at $-1$. In $\mathbb F^n$, part (iii) is the equivalence of (ii) and (iv) in \cref{thm:invertible}: $n$ columns are independent if and only if they span.

\subsection{Subspaces and the Dimension Formula}

\begin{proposition}[Dimension of a subspace]
\label{prop:subspace-dimension}
Let $V$ have dimension $n$, and let $U$ be a subspace of $V$. Then $U$ is finite-dimensional, $\dim U\le n$, and $\dim U=n$ if and only if $U=V$. Every basis of $U$ can be extended to a basis of $V$.
\end{proposition}

\begin{proof}
Every independent list in $U$ is independent in $V$ and has at most $n$ elements, by \cref{thm:basis-extension}\,(i). Choose an independent list $\mathbf u_1,\ldots,\mathbf u_k$ in $U$ of the largest possible length $k\le n$. It spans $U$: if some $\mathbf w\in U$ did not lie in its span, then $\mathbf u_1,\ldots,\mathbf u_k,\mathbf w$ would be a longer independent list in $U$, by \cref{lem:extend-independent}. So $U$ has a basis with $k\le n$ elements. If $k=n$, this basis consists of $n$ independent vectors of $V$ and is therefore a basis of $V$, by \cref{thm:basis-extension}\,(iii), so that $U=\operatorname{span}\{\mathbf u_1,\ldots,\mathbf u_n\}=V$. The last statement is \cref{thm:basis-extension}\,(ii).
\end{proof}

The proposition classifies the subspaces of $\mathbb R^3$ by their dimension: a subspace of dimension $0$ is $\{\mathbf 0\}$; one of dimension $1$ is spanned by a single nonzero vector and is a line through the origin; one of dimension $2$ is spanned by two vectors, neither a multiple of the other, and is a plane through the origin; and one of dimension $3$ is $\mathbb R^3$ itself. There are no others, and in particular there is no ``curved'' subspace.

The dimensions of two subspaces, of their intersection and of their sum are related in the same way as the numbers of elements of two finite sets, of their intersection and of their union (\cref{sec:counting-principles}).

\begin{theorem}[Dimension formula]
\label{thm:dimension-formula}
For finite-dimensional subspaces $U$ and $W$ of a vector space,
\[
\dim(U+W)+\dim(U\cap W)=\dim U+\dim W .
\]
In particular $\dim(U\oplus W)=\dim U+\dim W$.
\end{theorem}

\begin{proof}
The space $U+W$ is spanned by the union of a basis of $U$ and a basis of $W$, so it is finite-dimensional, and so is $U\cap W$, by \cref{prop:subspace-dimension}. Let $\mathbf z_1,\ldots,\mathbf z_k$ be a basis of $U\cap W$, and extend it to a basis $\mathbf z_1,\ldots,\mathbf z_k,\mathbf u_1,\ldots,\mathbf u_p$ of $U$ and to a basis $\mathbf z_1,\ldots,\mathbf z_k,\mathbf w_1,\ldots,\mathbf w_q$ of $W$ (\cref{prop:subspace-dimension}). We show that
\[
\mathbf z_1,\ldots,\mathbf z_k,\mathbf u_1,\ldots,\mathbf u_p,\mathbf w_1,\ldots,\mathbf w_q
\]
is a basis of $U+W$. It spans $U+W$, since every $\mathbf u\in U$ is a combination of the $\mathbf z_i$ and $\mathbf u_i$ and every $\mathbf w\in W$ of the $\mathbf z_i$ and $\mathbf w_i$. Suppose that
\[
\sum_ia_i\mathbf z_i+\sum_ib_i\mathbf u_i+\sum_ic_i\mathbf w_i=\mathbf 0 .
\]
Then $\sum_ic_i\mathbf w_i=-\sum_ia_i\mathbf z_i-\sum_ib_i\mathbf u_i$ lies in $W$ and in $U$, hence in $U\cap W$, and so $\sum_ic_i\mathbf w_i=\sum_id_i\mathbf z_i$ for some scalars $d_i$. Since the $\mathbf z_i$ and $\mathbf w_i$ together are independent, all $c_i$ and $d_i$ vanish. The relation reduces to $\sum_ia_i\mathbf z_i+\sum_ib_i\mathbf u_i=\mathbf 0$, and the independence of the basis of $U$ gives $a_i=b_i=0$. Counting, $\dim(U+W)=k+p+q=(k+p)+(k+q)-k=\dim U+\dim W-\dim(U\cap W)$.
\end{proof}

The formula has immediate geometric consequences. Two planes $U$ and $W$ through the origin of $\mathbb R^3$ satisfy $\dim(U+W)\le3$, so $\dim(U\cap W)\ge2+2-3=1$: they always share at least a line, as in panel~B of \cref{fig:subspaces}, and they share exactly a line unless they coincide. In $\mathbb R^4$, by contrast, two planes through the origin can meet only in $\mathbf 0$, for instance the planes spanned by $\mathbf e_1,\mathbf e_2$ and by $\mathbf e_3,\mathbf e_4$. In general, two subspaces of an $n$-dimensional space with $\dim U+\dim W>n$ have a nonzero vector in common. Conversely, if $U\cap W=\{\mathbf 0\}$ and $\dim U+\dim W=\dim V$, then $\dim(U\oplus W)=\dim V$, and $V=U\oplus W$ by \cref{prop:subspace-dimension}; this is the quickest way to prove that a space is a direct sum, and it shows again that $\mathbb R^3$ is the direct sum of the plane and the line of \cref{fig:subspaces}A.

\subsection{The Dimensions of the Examples}

\begin{proposition}[Dimensions of some spaces]
\label{prop:dimension-examples}
\leavevmode
\begin{enumerate}
\item The symmetric matrices in $\mathbb F^{n\times n}$ form a space of dimension $n(n+1)/2$, the antisymmetric ones a space of dimension $n(n-1)/2$, and the upper triangular ones a space of dimension $n(n+1)/2$.
\item The real sequences with $a_{k+2}=a_{k+1}+a_k$ for all $k\in\mathbb N$ form a space of dimension $2$.
\item The space $\mathcal P$ of all polynomials, the spaces $C[a,b]$ and $C^k[a,b]$ for $a<b$, and the space $\mathbb F^{\mathbb N}$ of all sequences are infinite-dimensional.
\end{enumerate}
\end{proposition}

\begin{proof}
(i) A symmetric matrix $A$ is determined by its entries on and above the diagonal, and $A=\sum_ia_{ii}E_{ii}+\sum_{i<j}a_{ij}(E_{ij}+E_{ji})$. So the $n$ matrices $E_{ii}$ and the $n(n-1)/2$ matrices $E_{ij}+E_{ji}$ with $i<j$ span the symmetric matrices, and they are independent, because a combination of them that vanishes has all its entries, which are the coefficients, equal to zero. Their number is $n+n(n-1)/2=n(n+1)/2$. An antisymmetric matrix has $a_{ii}=-a_{ii}$, hence $a_{ii}=0$, and in the same way the $n(n-1)/2$ matrices $E_{ij}-E_{ji}$ with $i<j$ form a basis. The matrix units $E_{ij}$ with $i\le j$ form a basis of the upper triangular matrices.

(ii) Let $\mathbf f=(1,0,1,1,2,3,\ldots)$ and $\mathbf g=(0,1,1,2,3,5,\ldots)$ be the solutions of the recurrence with the initial terms $(1,0)$ and $(0,1)$. A solution $\mathbf a$ and the solution $a_0\mathbf f+a_1\mathbf g$ agree in their first two terms, and then, by induction on $k$, in all terms, since the recurrence determines each term from the two before it; so $\mathbf f$ and $\mathbf g$ span the space. A relation $c_1\mathbf f+c_2\mathbf g=\mathbf 0$ gives $c_1=0$ and $c_2=0$ from the first two terms. So $(\mathbf f,\mathbf g)$ is a basis.

(iii) If $\mathcal P$ were finitely generated, of dimension $n$, then every independent list in it would have at most $n$ elements, by \cref{thm:basis-extension}\,(i); but $1,x,\ldots,x^n$ is an independent list with $n+1$ elements. The restrictions of $1,x,\ldots,x^n$ to $[a,b]$ are independent in $C^k[a,b]\subseteq C[a,b]$ as well, since a polynomial that vanishes on $[a,b]$ has infinitely many roots, and the same argument applies. In $\mathbb F^{\mathbb N}$, the sequences $\mathbf e_j$ with the term $1$ in position $j$ and zeros elsewhere are independent for every number of them, by looking at the terms in the positions $j$.
\end{proof}

The dimension of the antisymmetric matrices uses that $2a_{ii}=0$ implies $a_{ii}=0$, which fails over the field $\mathbb F_2$, where $2=0$ and every symmetric matrix is antisymmetric; this is the exception mentioned in the note of \cref{sec:vector-spaces}. Over $\mathbb R$ and $\mathbb C$ the dimensions $n(n+1)/2$ and $n(n-1)/2$ add up to $n^2=\dim\mathbb F^{n\times n}$, in accordance with the direct sum of \cref{sec:subspaces} and \cref{thm:dimension-formula}. The dimension of the space of solutions of the Fibonacci recurrence is the discrete counterpart of the dimension $2$ of the solutions of $f''=-\omega^2f$; in both cases the solutions are determined by two initial values, and Part~IV shows that the solutions of a linear recurrence or differential equation of order $k$ form a space of dimension $k$.

With dimensions available, the flats of \cref{sec:subspaces} acquire a dimension too: the \term{dimension of a flat} $\mathbf x_0+U$ is $\dim U$, which is well defined because a flat determines its direction (\cref{prop:flats}\,(ii)). A line is a flat of dimension $1$ and a plane one of dimension $2$, and the construction $\mathbf p+\operatorname{span}\{\mathbf v_1,\ldots,\mathbf v_k\}$ of \cref{sec:vectors-in-rn} with independent $\mathbf v_1,\ldots,\mathbf v_k$ gives a flat of dimension $k$, since these vectors form a basis of its direction. By \cref{thm:solution-structure} and the dimension of the null space computed in \cref{sec:four-subspaces}, the solution set of a consistent system $A\mathbf x=\mathbf b$ with $A\in\mathbb F^{m\times n}$ is a flat of dimension $n-\operatorname{rank}A$.

\begin{note}{Bases of infinite-dimensional spaces}
Every vector space has a basis in the sense of \cref{def:basis}, that is, a subset such that every vector is a combination of finitely many of its elements in exactly one way. For finitely generated spaces this is \cref{thm:dimension}; in general the proof uses the axiom of choice of the note in \cref{sec:sets}, usually in the form of Zorn's lemma, and the statement that every vector space has a basis is in fact equivalent to the axiom of choice~\cite{blass1984}. Such bases are called \emph{Hamel bases}, after Georg Hamel, who constructed one for the real numbers regarded as a vector space over $\mathbb Q$~\cite{hamel1905}. With a Hamel basis one can define additive functions $f(x+y)=f(x)+f(y)$ on $\mathbb R$ that are not of the form $cx$, by prescribing arbitrary values on the basis; these are the discontinuous additive functions mentioned in \cref{sec:general-powers}~\cite{aczel1966}. No Hamel basis of $\mathbb R$ over $\mathbb Q$ or of $C[a,b]$ over $\mathbb R$ can be written down explicitly, and they play no role in computation. The bases that matter in analysis are of a different kind: a Fourier series $\sum_k(a_k\cos kx+b_k\sin kx)$ expresses a function as an \emph{infinite} combination of the functions $\cos kx$ and $\sin kx$, which converges in a suitable sense. Such expansions require a notion of convergence in the vector space, and Volume~II develops them in the spaces $L^2$ of Lebesgue theory.
\end{note}

\begin{note}{Quotient spaces}
The flats with a fixed direction $U$ divide $V$ into disjoint pieces, by \cref{prop:flats}\,(iv), and they form a vector space of their own, the \emph{quotient space} $V/U$, with the operations $(\mathbf x+U)+(\mathbf y+U)\coloneqq(\mathbf x+\mathbf y)+U$ and $t(\mathbf x+U)\coloneqq t\mathbf x+U$, which do not depend on the chosen base points, by \cref{prop:flats}\,(i). Its zero vector is the flat $U$ itself. In the language of \cref{sec:equivalence}, $V/U$ is the set of equivalence classes of the relation $\mathbf x\sim\mathbf y\Leftrightarrow\mathbf x-\mathbf y\in U$. If $V$ is finite-dimensional, then $\dim V/U=\dim V-\dim U$: if $\mathbf u_1,\ldots,\mathbf u_k$ is a basis of $U$ extended to a basis $\mathbf u_1,\ldots,\mathbf u_k,\mathbf w_1,\ldots,\mathbf w_l$ of $V$, then the flats $\mathbf w_1+U,\ldots,\mathbf w_l+U$ form a basis of $V/U$, as one checks directly from the definitions. For the solution sets of the systems $A\mathbf x=\mathbf b$, $\mathbf b\in\mathcal C(A)$, which are the flats with direction $\mathcal N(A)$, the quotient $\mathbb F^n/\mathcal N(A)$ has the dimension $n-\dim\mathcal N(A)=\operatorname{rank}A$, the dimension of the space $\mathcal C(A)$ of right-hand sides; Chapter~16 explains this coincidence by a linear map between the two spaces. Quotients are used in algebra and analysis to identify objects that differ by something negligible, such as functions that agree almost everywhere in Volume~II; Axler treats them in detail~\cite[Sect.~3E]{axler2024}.
\end{note}

\section{Coordinates}
\label{sec:coordinates}

\subsection{Coordinate Vectors}

A basis turns vectors into lists of numbers. Once a basis of an $n$-dimensional space is fixed, every vector is described by the $n$ coefficients of its unique representation, and all computations in the space can be carried out with these numbers.

\begin{definition}[Coordinates]
\label{def:coordinates}
Let $\mathcal B=(\mathbf b_1,\ldots,\mathbf b_n)$ be a basis of $V$. For $\mathbf v\in V$, the scalars $t_1,\ldots,t_n$ with $\mathbf v=t_1\mathbf b_1+\cdots+t_n\mathbf b_n$ are the \term{coordinates} of $\mathbf v$ with respect to $\mathcal B$, and
\[
[\mathbf v]_{\mathcal B}\coloneqq(t_1,\ldots,t_n)\in\mathbb F^n
\]
is its \term{coordinate vector}.
\end{definition}

The coordinates are well defined by \cref{thm:basis-unique}, and they depend on the order of the basis vectors, which is why a basis is a list rather than a set. With respect to the standard basis of $\mathbb F^n$, the coordinate vector of $\mathbf x$ is $\mathbf x$ itself. With respect to the monomial basis $(1,x,\ldots,x^n)$ of $\mathcal P_n$, the coordinate vector of a polynomial is the list of its coefficients, and with respect to the basis $(\cos\omega x,\sin\omega x)$ of the solutions of $f''=-\omega^2f$, the coordinate vector of a solution $f$ is $(f(0),f'(0)/\omega)$, by \cref{lem:harmonic-unique}. In the space $\mathbb R_{>0}$ of \cref{sec:vector-spaces}, with the basis $(e)$, the coordinate of $x$ is $\ln x$: the logarithm is a coordinate map.

\begin{theorem}[The coordinate map]
\label{thm:coordinate-map}
Let $\mathcal B=(\mathbf b_1,\ldots,\mathbf b_n)$ be a basis of $V$.
\begin{enumerate}
\item The map $\mathbf v\mapsto[\mathbf v]_{\mathcal B}$ is linear: $[\mathbf u+\mathbf v]_{\mathcal B}=[\mathbf u]_{\mathcal B}+[\mathbf v]_{\mathcal B}$ and $[t\mathbf v]_{\mathcal B}=t[\mathbf v]_{\mathcal B}$.
\item It is a bijection from $V$ onto $\mathbb F^n$, whose inverse maps $(t_1,\ldots,t_n)$ to $t_1\mathbf b_1+\cdots+t_n\mathbf b_n$.
\item Vectors $\mathbf v_1,\ldots,\mathbf v_k\in V$ are linearly independent if and only if their coordinate vectors are, and they span $V$ if and only if their coordinate vectors span $\mathbb F^n$.
\end{enumerate}
\end{theorem}

\begin{proof}
(i) If $\mathbf u=\sum_js_j\mathbf b_j$ and $\mathbf v=\sum_jt_j\mathbf b_j$, then $\mathbf u+\mathbf v=\sum_j(s_j+t_j)\mathbf b_j$ and $t\mathbf v=\sum_j(tt_j)\mathbf b_j$, and by uniqueness these are the representations of $\mathbf u+\mathbf v$ and $t\mathbf v$. (ii) The map $g\colon\mathbb F^n\to V$, $\mathbf t\mapsto\sum_jt_j\mathbf b_j$, satisfies $[g(\mathbf t)]_{\mathcal B}=\mathbf t$ by the definition of coordinates, and $g([\mathbf v]_{\mathcal B})=\mathbf v$ for every $\mathbf v$; so the two maps are inverse to each other (\cref{sec:injections}). (iii) By (i), $[\sum_ic_i\mathbf v_i]_{\mathcal B}=\sum_ic_i[\mathbf v_i]_{\mathcal B}$, and since the coordinate map is a bijection that maps $\mathbf 0$ to $\mathbf 0$, a combination of the $\mathbf v_i$ is zero if and only if the same combination of their coordinate vectors is zero. This gives the statement about independence. Similarly, a vector $\mathbf v$ is a combination of the $\mathbf v_i$ with coefficients $c_i$ if and only if $[\mathbf v]_{\mathcal B}$ is the combination of the $[\mathbf v_i]_{\mathcal B}$ with the same coefficients, and every vector of $\mathbb F^n$ is a coordinate vector.
\end{proof}

The theorem says that an $n$-dimensional vector space over $\mathbb F$, whatever its elements are, behaves exactly like $\mathbb F^n$ once a basis is chosen: sums correspond to sums, multiples to multiples, and independence and spanning can be tested on the coordinate vectors by elimination (\cref{prop:independence-columns}). In the language of Chapter~16, the coordinate map is an \emph{isomorphism}, and any two vector spaces of the same finite dimension over the same field are isomorphic. This is why the theory of finite-dimensional spaces is, at bottom, the theory of $\mathbb F^n$ and of matrices; but the identification depends on the choice of a basis, and much of linear algebra consists of choosing the basis well.

For example, are the polynomials $q_1=1+x$, $q_2=1-x$ and $q_3=x+x^2$ a basis of $\mathcal P_2$, and what are the coordinates of $p=1+x+x^2$ with respect to them? Their coordinate vectors with respect to the monomial basis are $(1,1,0)$, $(1,-1,0)$ and $(0,1,1)$, and we must decide whether the matrix
\[
Q=\begin{bmatrix}1&1&0\\1&-1&1\\0&0&1\end{bmatrix}
\]
with these columns has rank $3$. Subtracting the first row from the second gives an upper triangular matrix with the pivots $1$, $-2$ and $1$, so the rank is $3$, the three coordinate vectors are independent, and by \cref{thm:coordinate-map}\,(iii) and \cref{thm:basis-extension}\,(iii) the three polynomials form a basis of $\mathcal P_2$. The coordinates of $p$ solve $Q\mathbf t=(1,1,1)$, the coordinate vector of $p$; back substitution gives $t_3=1$, then $t_1-t_2=0$ and $t_1+t_2=1$, so $\mathbf t=(1/2,1/2,1)$, and indeed $(1+x)/2+(1-x)/2+(x+x^2)=1+x+x^2$.

\subsection{Change of Basis}

A vector has different coordinates with respect to different bases, and the coordinates with respect to one basis are computed from those with respect to another by a matrix.

\begin{theorem}[Change of basis]
\label{thm:change-of-basis}
Let $\mathcal B=(\mathbf b_1,\ldots,\mathbf b_n)$ and $\mathcal B'=(\mathbf b_1',\ldots,\mathbf b_n')$ be bases of $V$, and let $S\in\mathbb F^{n\times n}$ be the matrix whose $j$th column is $[\mathbf b_j']_{\mathcal B}$. Then
\[
[\mathbf v]_{\mathcal B}=S\,[\mathbf v]_{\mathcal B'}\qquad\text{for every }\mathbf v\in V .
\]
The matrix $S$ is invertible, and $S^{-1}$ is the matrix whose $j$th column is $[\mathbf b_j]_{\mathcal B'}$, so that $[\mathbf v]_{\mathcal B'}=S^{-1}[\mathbf v]_{\mathcal B}$.
\end{theorem}

\begin{proof}
If $[\mathbf v]_{\mathcal B'}=\mathbf t$, that is, $\mathbf v=\sum_jt_j\mathbf b_j'$, then by the linearity of the coordinate map
\[
[\mathbf v]_{\mathcal B}=\sum_jt_j[\mathbf b_j']_{\mathcal B}=S\mathbf t,
\]
the combination of the columns of $S$ with the coefficients $t_j$. Let $T$ be the matrix whose $j$th column is $[\mathbf b_j]_{\mathcal B'}$; by the same argument with the roles of the bases exchanged, $[\mathbf v]_{\mathcal B'}=T[\mathbf v]_{\mathcal B}$ for every $\mathbf v$. Then $ST[\mathbf v]_{\mathcal B}=[\mathbf v]_{\mathcal B}$ for every $\mathbf v$, and taking $\mathbf v=\mathbf b_j$, whose coordinate vector is $\mathbf e_j$, shows that the $j$th column of $ST$ is $\mathbf e_j$. So $ST=I$, and by \cref{thm:invertible}, $S$ is invertible with $S^{-1}=T$.
\end{proof}

The matrix $S$ is the \term{change-of-basis matrix} from $\mathcal B'$ to $\mathcal B$: it converts coordinates with respect to $\mathcal B'$ into coordinates with respect to $\mathcal B$, and its columns are the new basis vectors expressed in the old basis. The direction is easy to confuse, and the safest way to remember it is the formula $[\mathbf v]_{\mathcal B}=S[\mathbf v]_{\mathcal B'}$ applied to $\mathbf v=\mathbf b_j'$, which gives back the definition of the $j$th column. Conversely, every invertible matrix $S$ is a change-of-basis matrix: if $\mathcal B$ is a basis and $S$ is invertible, the vectors $\mathbf b_j'$ with $[\mathbf b_j']_{\mathcal B}=\mathbf s_j$, the $j$th column of $S$, form a basis, since their coordinate vectors are $n$ independent columns.

The most important case is $V=\mathbb F^n$ with $\mathcal B$ the standard basis. Then $[\mathbf b_j']_{\mathcal B}=\mathbf b_j'$, so $S$ is simply the matrix with the columns $\mathbf b_1',\ldots,\mathbf b_n'$, and the coordinates $\mathbf t=[\mathbf v]_{\mathcal B'}$ of a vector $\mathbf v$ with respect to the new basis are the solution of the linear system $S\mathbf t=\mathbf v$. As \cref{sec:invertible} advised, one solves this system rather than forming $S^{-1}$. For the basis $\mathbf b_1'=(1,1)$, $\mathbf b_2'=(-1,2)$ of $\mathbb R^2$ and $\mathbf v=(1,4)$,
\[
S=\begin{bmatrix}1&-1\\1&2\end{bmatrix},\qquad S\mathbf t=\begin{bmatrix}1\\4\end{bmatrix}\quad\Longrightarrow\quad\mathbf t=\begin{bmatrix}2\\1\end{bmatrix},
\]
so that $\mathbf v=2\mathbf b_1'+\mathbf b_2'$. \Cref{fig:coordinates} shows the two coordinate grids: the same point has the coordinates $(1,4)$ in the standard grid and $(2,1)$ in the slanted grid of the new basis, the construction of \cref{fig:linear-combinations} read in the other direction.

\begin{figure}
\sdpanel{A}{\begin{tikzpicture}[scale=0.62, font=\scriptsize]
\begin{scope}
\clip (-2.4,-0.8) rectangle (2.9,5.0);
\foreach \s in {-3,...,3} {\draw[black!25, line width=0.35pt] (\s,-1) -- (\s,5.2);}
\foreach \t in {-1,...,5} {\draw[black!25, line width=0.35pt] (-3,\t) -- (3,\t);}
\end{scope}
\draw[line width=0.4pt, -{Latex[length=1.5mm]}] (-2.4,0) -- (3.1,0) node[below] {$x_1$};
\draw[line width=0.4pt, -{Latex[length=1.5mm]}] (0,-0.8) -- (0,5.2) node[left] {$x_2$};
\draw[black!60, dashed, line width=0.5pt] (1,0) -- (1,4);
\draw[line width=0.8pt, -{Latex[length=2mm]}] (0,0) -- (1,0) node[below] {$\mathbf e_1$};
\draw[line width=0.8pt, -{Latex[length=2mm]}] (0,0) -- (0,1) node[left] {$\mathbf e_2$};
\fill (1,4) circle (2.4pt);
\node[right] at (1.05,4.05) {$\mathbf v=\mathbf e_1+4\mathbf e_2$};
\end{tikzpicture}}\hfill
\sdpanel{B}{\begin{tikzpicture}[scale=0.62, font=\scriptsize]
\begin{scope}
\clip (-2.4,-0.8) rectangle (2.9,5.0);
\foreach \s in {-4,...,4} {\draw[black!25, line width=0.35pt] ({\s-4},{\s-8}) -- ({\s+4},{\s+8});}
\foreach \t in {-4,...,4} {\draw[black!25, line width=0.35pt] ({-\t-6},{2*\t-6}) -- ({-\t+6},{2*\t+6});}
\end{scope}
\draw[line width=0.4pt, -{Latex[length=1.5mm]}] (-2.4,0) -- (3.1,0) node[below] {$x_1$};
\draw[line width=0.4pt, -{Latex[length=1.5mm]}] (0,-0.8) -- (0,5.2) node[left] {$x_2$};
\draw[black!60, dashed, line width=0.5pt] (2,2) -- (1,4);
\draw[line width=0.8pt, -{Latex[length=2mm]}] (0,0) -- (1,1) node[right] {$\mathbf b_1'$};
\draw[line width=0.8pt, -{Latex[length=2mm]}] (0,0) -- (-1,2) node[left] {$\mathbf b_2'$};
\draw[line width=1.1pt, -{Latex[length=2mm]}] (1,1) -- (2,2);
\fill (1,4) circle (2.4pt);
\node[right] at (1.05,4.05) {$\mathbf v=2\mathbf b_1'+\mathbf b_2'$};
\end{tikzpicture}}
\caption{One vector, two coordinate systems. (A) With respect to the standard basis, $\mathbf v=(1,4)$ has the coordinates $1$ and $4$. (B) With respect to the basis $\mathbf b_1'=(1,1)$, $\mathbf b_2'=(-1,2)$, whose integer combinations form the slanted grid, the same vector has the coordinates $2$ and $1$: two steps along $\mathbf b_1'$ and one step along $\mathbf b_2'$ (dashed). The coordinates are found by solving $S\mathbf t=\mathbf v$ with the matrix $S$ whose columns are $\mathbf b_1'$ and $\mathbf b_2'$.}
\label{fig:coordinates}
\end{figure}

\subsection{Three Bases of \texorpdfstring{$\mathcal P_2$}{P2}}

The space of polynomials of degree at most $2$ has several useful bases, and the change of basis between them connects the theory of this chapter with the interpolation problem of \cref{sec:special-systems}. Fix the three points $0$, $1$, $2$. Besides the monomial basis $\mathcal B_{\mathrm M}=(1,x,x^2)$ we consider the \term{Lagrange basis} $\mathcal B_{\mathrm L}=(L_0,L_1,L_2)$ and the \term{Newton basis} $\mathcal B_{\mathrm N}=(1,x,x(x-1))$, where
\[
L_0(x)=\frac{(x-1)(x-2)}2,\qquad L_1(x)=-x(x-2),\qquad L_2(x)=\frac{x(x-1)}2 .
\]
Each $L_i$ has degree $2$, equals $1$ at the point $i$ and vanishes at the other two points. For every $p\in\mathcal P_2$, the polynomial $q\coloneqq p(0)L_0+p(1)L_1+p(2)L_2$ takes the values $p(0)$, $p(1)$, $p(2)$ at the three points, so $p-q\in\mathcal P_2$ has three roots and is zero (\cref{cor:factor}). So the $L_i$ span $\mathcal P_2$, hence form a basis, by \cref{thm:basis-extension}\,(iii), and
\[
[p]_{\mathcal B_{\mathrm L}}=\bigl(p(0),p(1),p(2)\bigr):
\]
the coordinates with respect to the Lagrange basis are the values at the points. The change-of-basis matrix from $\mathcal B_{\mathrm M}$ to $\mathcal B_{\mathrm L}$ has as its $j$th column the values of the $j$th monomial at the points, and it is the Vandermonde matrix $V$ of \cref{sec:special-systems}; the change-of-basis matrix in the other direction has as columns the coefficient vectors of the $L_i$:
\[
[p]_{\mathcal B_{\mathrm L}}=V[p]_{\mathcal B_{\mathrm M}},\quad V=\begin{bmatrix}1&0&0\\1&1&1\\1&2&4\end{bmatrix};\qquad
[p]_{\mathcal B_{\mathrm M}}=V^{-1}[p]_{\mathcal B_{\mathrm L}},\quad V^{-1}=\begin{bmatrix}1&0&0\\-\frac32&2&-\frac12\\\frac12&-1&\frac12\end{bmatrix}.
\]
Interpolation, finding the coefficients of the polynomial with given values, is therefore a change of basis from the Lagrange basis to the monomial basis, and \cref{thm:interpolation}, the invertibility of $V$, is the statement that the Lagrange polynomials form a basis. For the values $1,3,7$ of \cref{sec:special-systems}, $V^{-1}(1,3,7)=(1,1,1)$, the coefficients of $1+x+x^2$.

The Newton basis lies between the two. The values of its elements at the points $0,1,2$ are the columns of a lower triangular matrix, so the change of basis from $\mathcal B_{\mathrm N}$ to $\mathcal B_{\mathrm L}$ is
\[
[p]_{\mathcal B_{\mathrm L}}=W[p]_{\mathcal B_{\mathrm N}},\qquad W=\begin{bmatrix}1&0&0\\1&1&0\\1&2&2\end{bmatrix},
\]
and the Newton coordinates of the polynomial with the values $1,3,7$ are found by forward substitution: $c_0=1$, $c_0+c_1=3$ gives $c_1=2$, and $c_0+2c_1+2c_2=7$ gives $c_2=1$; indeed $1+2x+x(x-1)=1+x+x^2$. The triangular shape means that a new point can be added without changing the coordinates already computed, which is the reason for the popularity of the Newton form; with $d+1$ points, the Newton coordinates cost $O(d^2)$ operations against the $O(d^3)$ of a general system. Volume~XI develops both forms and the choice of the points.

\section{The Four Fundamental Subspaces}
\label{sec:four-subspaces}

\subsection{Four Subspaces of a Matrix}

Every matrix $A\in\mathbb F^{m\times n}$ comes with two subspaces met in \cref{sec:subspaces}, the null space $\mathcal N(A)\subseteq\mathbb F^n$ and the column space $\mathcal C(A)\subseteq\mathbb F^m$, and the transpose $A^{\mathsf T}\in\mathbb F^{n\times m}$ comes with two more.

\begin{definition}[The four fundamental subspaces]
\label{def:four-subspaces}
For $A\in\mathbb F^{m\times n}$, the four \term{fundamental subspaces} are:
\begin{enumerate}
\item the \term{column space} $\mathcal C(A)=\{A\mathbf x: \mathbf x\in\mathbb F^n\}\subseteq\mathbb F^m$, the span of the columns of $A$;
\item the \term{null space} $\mathcal N(A)=\{\mathbf x\in\mathbb F^n: A\mathbf x=\mathbf 0\}\subseteq\mathbb F^n$;
\item the \term{row space} $\mathcal C(A^{\mathsf T})\subseteq\mathbb F^n$, the span of the rows of $A$, written as columns;
\item the \term{left null space} $\mathcal N(A^{\mathsf T})=\{\mathbf y\in\mathbb F^m: A^{\mathsf T}\mathbf y=\mathbf 0\}\subseteq\mathbb F^m$.
\end{enumerate}
\end{definition}

The name of the last space comes from transposing its defining equation: $A^{\mathsf T}\mathbf y=\mathbf 0$ is the same as $\mathbf y^{\mathsf T}A=\mathbf 0^{\mathsf T}$, so the left null space consists of the vectors that annihilate $A$ from the left. Such a vector records a linear relation among the rows of $A$, since $\mathbf y^{\mathsf T}A=\sum_iy_i\mathbf r_i^{\mathsf T}$ is a combination of the rows $\mathbf r_i^{\mathsf T}$ (\cref{thm:four-readings}\,(iii)); \cref{prop:singular-image} produced such a vector for every singular square matrix. Two of the spaces live in $\mathbb F^n$, the space of the unknowns, and two in $\mathbb F^m$, the space of the right-hand sides. Elimination gives bases of all four at once. The first step is to see what row operations do to them.

\begin{lemma}[Row operations and the four subspaces]
\label{lem:row-operations-subspaces}
Let $A\in\mathbb F^{m\times n}$, let $E\in\mathbb F^{m\times m}$ be invertible, and let $B\coloneqq EA$, with the columns $\mathbf a_j$ and $\mathbf b_j$.
\begin{enumerate}
\item $\mathcal N(B)=\mathcal N(A)$.
\item $\mathcal C(B^{\mathsf T})=\mathcal C(A^{\mathsf T})$: row operations do not change the row space.
\item The columns of $A$ and of $B$ satisfy the same linear relations: $\sum_jx_j\mathbf b_j=\mathbf 0$ if and only if $\sum_jx_j\mathbf a_j=\mathbf 0$. In particular, a set of columns of $A$ is linearly independent if and only if the corresponding columns of $B$ are, and $\mathbf a_k=\sum_{j\in J}c_j\mathbf a_j$ if and only if $\mathbf b_k=\sum_{j\in J}c_j\mathbf b_j$.
\end{enumerate}
\end{lemma}

\begin{proof}
(i) is \cref{thm:row-equivalence}: $EA\mathbf x=\mathbf 0$ if and only if $A\mathbf x=\mathbf 0$, since $E$ is invertible. (ii) By the row reading of the product (\cref{thm:four-readings}\,(iii)), the $i$th row of $EA$ is $\sum_ke_{ik}\mathbf r_k^{\mathsf T}$, a combination of the rows of $A$. So every row of $B$ lies in the row space of $A$, and by \cref{prop:span-smallest}, $\mathcal C(B^{\mathsf T})\subseteq\mathcal C(A^{\mathsf T})$. Since $A=E^{-1}B$, the reverse inclusion holds as well. (iii) $\sum_jx_j\mathbf b_j=B\mathbf x$ and $\sum_jx_j\mathbf a_j=A\mathbf x$, so this is (i) once more. A relation among some of the columns is a relation among all of them with zero coefficients for the others, and $\mathbf a_k=\sum_{j\in J}c_j\mathbf a_j$ is a relation among the columns.
\end{proof}

The column space itself does change under row operations. The matrix with the single column $(1,2)$ reduces to the column $(1,0)$, and the column spaces are two different lines of $\mathbb R^2$. What is preserved is the pattern of dependence among the columns, and this is all we need.

\subsection{Bases from the Reduced Row Echelon Form}

Let $A\in\mathbb F^{m\times n}$ have the reduced row echelon form $R=GA$, where $G\in\mathbb F^{m\times m}$ is a product of elementary matrices and hence invertible, let $r$ be the number of pivots, $p_1<\cdots<p_r$ the pivot columns and $f_1<\cdots<f_{n-r}$ the free columns, and let $\mathbf s_1,\ldots,\mathbf s_{n-r}$ be the special solutions of \cref{thm:homogeneous-solutions}.

\begin{theorem}[Bases of the four subspaces]
\label{thm:four-bases}
With this notation:
\begin{enumerate}
\item the pivot columns $\mathbf a_{p_1},\ldots,\mathbf a_{p_r}$ \emph{of $A$} form a basis of $\mathcal C(A)$; more precisely, they are independent, and every free column of $A$ is the combination
\[
\mathbf a_{f_k}=\sum_{i=1}^rr_{if_k}\,\mathbf a_{p_i}\qquad(k=1,\ldots,n-r)
\]
with the entries of the corresponding column of $R$ as coefficients;
\item the $r$ nonzero rows of $R$ form a basis of the row space $\mathcal C(A^{\mathsf T})$;
\item the special solutions $\mathbf s_1,\ldots,\mathbf s_{n-r}$ form a basis of $\mathcal N(A)$;
\item the last $m-r$ rows of $G$ form a basis of the left null space $\mathcal N(A^{\mathsf T})$.
\end{enumerate}
\end{theorem}

\begin{proof}
(i) In $R$, the pivot column $p_i$ is the standard basis vector $\mathbf e_i\in\mathbb F^m$, by (iii) and (iv) of \cref{def:echelon}, so the pivot columns of $R$ are independent. A free column $f_k$ of $R$ has the entries $r_{1f_k},\ldots,r_{rf_k}$ and zeros below the $r$th row, since these rows of $R$ vanish; so it equals $\sum_{i\le r}r_{if_k}\mathbf e_i$, the combination of the pivot columns with these coefficients. By \cref{lem:row-operations-subspaces}\,(iii), the pivot columns of $A$ are independent, and each free column of $A$ is the same combination of the pivot columns of $A$. So the pivot columns span all columns of $A$, hence $\mathcal C(A)$, and they form a basis.

(ii) The row space of $A$ equals that of $R$, by \cref{lem:row-operations-subspaces}\,(ii), and the zero rows of $R$ contribute nothing to it, so the $r$ nonzero rows span it. They are independent: in a combination $\sum_{i\le r}c_i(\text{row }i\text{ of }R)$, the entry in the pivot column $p_k$ is $c_k$, since the pivot column $p_k$ of $R$ is $\mathbf e_k$; so the combination vanishes only if all $c_k$ vanish.

(iii) By \cref{thm:homogeneous-solutions}, every solution of $A\mathbf x=\mathbf 0$ is a combination $\sum_kt_k\mathbf s_k$ of the special solutions, so they span $\mathcal N(A)$, and the coefficients of a solution $\mathbf x$ are determined by it, $t_k=x_{f_k}$. In particular the zero vector has only the representation with zero coefficients, and the special solutions are independent.

(iv) Let $\mathbf y\in\mathbb F^m$, and put $\mathbf z\coloneqq(G^{-1})^{\mathsf T}\mathbf y$, so that $\mathbf y^{\mathsf T}=\mathbf z^{\mathsf T}G$ and $\mathbf y^{\mathsf T}A=\mathbf z^{\mathsf T}GA=\mathbf z^{\mathsf T}R$. Now $\mathbf z^{\mathsf T}R=\sum_iz_i(\text{row }i\text{ of }R)$, and since the first $r$ rows of $R$ are independent by (ii) and the others are zero, $\mathbf z^{\mathsf T}R=\mathbf 0^{\mathsf T}$ if and only if $z_1=\cdots=z_r=0$. Hence $\mathbf y\in\mathcal N(A^{\mathsf T})$ if and only if $\mathbf y^{\mathsf T}=\sum_{i>r}z_i\,\mathbf e_i^{\mathsf T}G$ for some $z_{r+1},\ldots,z_m$, that is, if and only if $\mathbf y^{\mathsf T}$ is a combination of the rows $r+1,\ldots,m$ of $G$. These rows span $\mathcal N(A^{\mathsf T})$, and they are independent, because $G$ is invertible: a relation $\mathbf c^{\mathsf T}G=\mathbf 0^{\mathsf T}$ among its rows gives $\mathbf c^{\mathsf T}=\mathbf 0^{\mathsf T}G^{-1}=\mathbf 0^{\mathsf T}$.
\end{proof}

Part (i) deserves emphasis: the basis of the column space consists of the pivot columns \emph{of $A$}, not of $R$, whose column space is in general a different subspace. The reduced form only tells us \emph{which} columns to take, and it tells us, in its free columns, how the remaining columns are combined from them. Part (iii) is what \cref{ch:linear-systems} anticipated after \cref{thm:homogeneous-solutions}, and part (iv) generalises \cref{prop:singular-image}, whose vector $\mathbf y$ is the transposed last row of $G$.

\subsection{The Dimensions and the Rank of the Transpose}

Counting the bases of \cref{thm:four-bases} gives the dimensions of the four subspaces, and with them the theorem of \cref{ch:linear-systems} that was proved there only in outline.

\begin{theorem}[Dimensions of the four subspaces]
\label{thm:four-dimensions}
Let $A\in\mathbb F^{m\times n}$ have rank $r$. Then
\[
\dim\mathcal C(A)=\dim\mathcal C(A^{\mathsf T})=r,\qquad\dim\mathcal N(A)=n-r,\qquad\dim\mathcal N(A^{\mathsf T})=m-r .
\]
In particular, $\operatorname{rank}A^{\mathsf T}=\operatorname{rank}A$, which is \cref{thm:rank-transpose}, and
\[
\dim\mathcal C(A)+\dim\mathcal N(A)=n,\qquad\dim\mathcal C(A^{\mathsf T})+\dim\mathcal N(A^{\mathsf T})=m .
\]
\end{theorem}

\begin{proof}
The four dimensions are the numbers of vectors in the bases of \cref{thm:four-bases}: $r$ pivot columns, $r$ nonzero rows of $R$, $n-r$ special solutions and $m-r$ rows of $G$. For the transpose, apply what has been shown to $A^{\mathsf T}\in\mathbb F^{n\times m}$: its row space has dimension $\operatorname{rank}A^{\mathsf T}$. The rows of $A^{\mathsf T}$ are the columns of $A$, so the row space of $A^{\mathsf T}$ is $\mathcal C(A)$, whose dimension is $r=\operatorname{rank}A$. Hence $\operatorname{rank}A^{\mathsf T}=\operatorname{rank}A$. The last two equations are $r+(n-r)=n$ and $r+(m-r)=m$.
\end{proof}

This completes the proof of \cref{thm:rank-transpose}, and we check that it is not circular. It rests on \cref{thm:four-bases}, which uses the reduced row echelon form, \cref{thm:row-equivalence,thm:homogeneous-solutions} and the definition of independence, and on the fact that the dimension is well defined (\cref{thm:dimension}), which rests on \cref{thm:more-unknowns}; none of these depends on \cref{thm:rank-transpose}. In particular we have not used \cref{prop:rank-product}, whose second inequality \cref{ch:linear-systems} derived from \cref{thm:rank-transpose}; that derivation is now complete as well. The essential point of the proof is that the rank, defined by counting pivots, is the dimension of \emph{two} different spaces, the span of the columns and the span of the rows, and the dimension of a space does not depend on how it is computed.

The first of the two equations at the end of the theorem is the \term{rank--nullity theorem} in its matrix form: the rank of $A$ and the dimension of its null space, its \term{nullity}, add up to the number of columns. It is the precise form of the counting of \cref{ch:linear-systems}: each of the $n$ unknowns is either a pivot variable, determined by the equations, or a free variable, contributing one dimension to the solutions. Chapter~16 states the theorem for linear maps between arbitrary finite-dimensional spaces, where it reads $\dim(\text{image})+\dim(\text{kernel})=\dim(\text{domain})$. The theorem also gives the dimension of the flats of \cref{sec:basis-dimension}: a hyperplane $\{\mathbf x\in\mathbb F^n: \mathbf a^{\mathsf T}\mathbf x=c\}$ with $\mathbf a\ne\mathbf 0$ is nonempty, since the $1\times n$ matrix $\mathbf a^{\mathsf T}$ has rank $1$ and \cref{thm:rank-solutions}\,(iii) applies, and its direction $\mathcal N(\mathbf a^{\mathsf T})$ has dimension $n-1$.

\subsection{A Complete Example}

We compute the four subspaces of the matrix of \cref{sec:echelon},
\[
A=\begin{bmatrix}1&2&0&3\\2&4&1&8\\-1&-2&1&-1\end{bmatrix}\in\mathbb R^{3\times4}.
\]
To obtain $G$ along with $R$, we apply the row operations to the augmented matrix $[A\mid I]$: subtracting twice the first row from the second, adding the first row to the third, and subtracting the new second row from the new third give
\[
[A\mid I]\;\longrightarrow\;
\left[\begin{array}{rrrr|rrr}1&2&0&3&1&0&0\\0&0&1&2&-2&1&0\\0&0&1&2&1&0&1\end{array}\right]
\;\longrightarrow\;
\left[\begin{array}{rrrr|rrr}1&2&0&3&1&0&0\\0&0&1&2&-2&1&0\\0&0&0&0&3&-1&1\end{array}\right]=[R\mid G].
\]
The left part $R$ is already in reduced row echelon form, with the pivot columns $1$ and $3$, the free columns $2$ and $4$, and $r=2$, and the right part is the matrix $G$ with $GA=R$. \Cref{thm:four-bases} gives the following bases.
\begin{itemize}
\item Column space $\mathcal C(A)\subseteq\mathbb R^3$, dimension $2$: the pivot columns $(1,2,-1)$ and $(0,1,1)$ of $A$. The free columns are $\mathbf a_2=2\mathbf a_1$ and $\mathbf a_4=3\mathbf a_1+2\mathbf a_3$, with the coefficients read off from the columns $2$ and $4$ of $R$; for instance $3(1,2,-1)+2(0,1,1)=(3,8,-1)$.
\item Row space $\mathcal C(A^{\mathsf T})\subseteq\mathbb R^4$, dimension $2$: the rows $(1,2,0,3)$ and $(0,0,1,2)$ of $R$.
\item Null space $\mathcal N(A)\subseteq\mathbb R^4$, dimension $4-2=2$: the special solutions $\mathbf s_1=(-2,1,0,0)$ and $\mathbf s_2=(-3,0,-2,1)$.
\item Left null space $\mathcal N(A^{\mathsf T})\subseteq\mathbb R^3$, dimension $3-2=1$: the last row $(3,-1,1)$ of $G$. Indeed $3\mathbf r_1-\mathbf r_2+\mathbf r_3=\mathbf 0$ for the rows of $A$, the relation noticed in \cref{sec:echelon}.
\end{itemize}
The dimensions add up as \cref{thm:four-dimensions} requires: $2+2=4=n$ and $2+1=3=m$. The example also illustrates the use of the left null space: the system $A\mathbf x=\mathbf b$ is consistent exactly when $3b_1-b_2+b_3=0$, the condition found in \cref{sec:echelon}; \cref{sec:fredholm} proves that the left null space always describes the consistent right-hand sides in this way. One more observation is worth making. Each basis vector of the row space is orthogonal to each special solution, for instance $\langle(1,2,0,3),(-3,0,-2,1)\rangle=-3+3=0$, and the vector $(3,-1,1)$ is orthogonal to both basis vectors of the column space. This is no accident, as the following proposition shows.

\begin{proposition}[Row space and null space of a real matrix]
\label{prop:row-null-real}
Let $A\in\mathbb R^{m\times n}$ have rank $r$.
\begin{enumerate}
\item Every vector of $\mathcal C(A^{\mathsf T})$ is orthogonal to every vector of $\mathcal N(A)$, and $\mathbb R^n=\mathcal C(A^{\mathsf T})\oplus\mathcal N(A)$.
\item Every vector of $\mathcal C(A)$ is orthogonal to every vector of $\mathcal N(A^{\mathsf T})$, and $\mathbb R^m=\mathcal C(A)\oplus\mathcal N(A^{\mathsf T})$.
\item The map $\mathbf x\mapsto A\mathbf x$ maps $\mathcal C(A^{\mathsf T})$ bijectively onto $\mathcal C(A)$, and every $\mathbf x\in\mathbb R^n$ has $A\mathbf x=A\mathbf x_{\mathrm r}$, where $\mathbf x=\mathbf x_{\mathrm r}+\mathbf x_{\mathrm n}$ is its decomposition with $\mathbf x_{\mathrm r}\in\mathcal C(A^{\mathsf T})$ and $\mathbf x_{\mathrm n}\in\mathcal N(A)$.
\end{enumerate}
\end{proposition}

\begin{proof}
(i) A vector of the row space has the form $A^{\mathsf T}\mathbf y$, and for $\mathbf x\in\mathcal N(A)$, $\langle A^{\mathsf T}\mathbf y,\mathbf x\rangle=\mathbf y^{\mathsf T}A\mathbf x=0$ (\cref{prop:transpose-rules}). If $\mathbf x$ lies in both spaces, it is orthogonal to itself, $\|\mathbf x\|^2=\langle\mathbf x,\mathbf x\rangle=0$, and so $\mathbf x=\mathbf 0$. The sum is therefore direct, of dimension $r+(n-r)=n$ by \cref{thm:dimension-formula,thm:four-dimensions}, and it is all of $\mathbb R^n$ by \cref{prop:subspace-dimension}. (ii) is (i) for $A^{\mathsf T}$, whose row space is $\mathcal C(A)$. (iii) $A\mathbf x=A\mathbf x_{\mathrm r}+A\mathbf x_{\mathrm n}=A\mathbf x_{\mathrm r}$. So every vector $A\mathbf x$ of the column space is the image of a vector of the row space, and if two vectors $\mathbf u,\mathbf w$ of the row space have the same image, then $\mathbf u-\mathbf w$ lies in the row space and in the null space, and is zero by (i).
\end{proof}

The proposition uses the inner product, and over $\mathbb C$ it is false as stated: for the $1\times2$ matrix $A$ with the entries $1$ and $i$, the vector $\mathbf x=(1,i)$ spans the row space $\mathcal C(A^{\mathsf T})$, and $A\mathbf x=1+i^2=0$, so $\mathbf x$ lies in the null space as well. The cure is to replace the transpose by the conjugate transpose of \cref{sec:complex-matrices}, which is natural for the Hermitian inner product; with $\mathcal C(A^{\mathsf H})$ and $\mathcal N(A^{\mathsf H})$ in place of $\mathcal C(A^{\mathsf T})$ and $\mathcal N(A^{\mathsf T})$, the proposition holds for complex matrices with the same proof. Over a general field there is no inner product at all, and in particular none with $\langle\mathbf x,\mathbf x\rangle=0$ only for $\mathbf x=\mathbf0$. Chapter~18 shows that the row space consists of \emph{all} vectors orthogonal to the null space, its orthogonal complement, and builds the theory of orthogonal projections on this fact. \Cref{fig:four-subspaces} summarises the four subspaces and the action of $A$ on them.

\begin{figure}
\centering
\begin{tikzpicture}[font=\scriptsize]
\fill[sdfill] (0.2,0) rectangle (2.4,2.2);
\draw[line width=0.5pt] (0.2,0) rectangle (2.4,2.2);
\fill[black!5] (2.4,0) rectangle (4.2,-1.5);
\draw[line width=0.5pt] (2.4,0) rectangle (4.2,-1.5);
\fill[sdfill] (7.4,0) rectangle (9.6,2.2);
\draw[line width=0.5pt] (7.4,0) rectangle (9.6,2.2);
\fill[black!5] (5.6,0) rectangle (7.4,-1.5);
\draw[line width=0.5pt] (5.6,0) rectangle (7.4,-1.5);
\fill (2.4,0) circle (1.6pt); \node[above right] at (2.4,0) {$\mathbf 0$};
\fill (7.4,0) circle (1.6pt); \node[below right] at (7.4,0) {$\mathbf 0$};
\node[align=center] at (1.3,1.45) {row space\\$\mathcal C(A^{\mathsf T})$\\dimension $r$};
\node[align=center] at (3.3,-0.62) {null space\\$\mathcal N(A)$\\dim.\ $n-r$};
\node[align=center] at (8.5,1.45) {column space\\$\mathcal C(A)$\\dimension $r$};
\node[align=center] at (6.5,-0.68) {left null space\\$\mathcal N(A^{\mathsf T})$\\dim.\ $m-r$};
\fill (1.6,0.5) circle (1.4pt); \node[left] at (1.6,0.5) {$\mathbf x_{\mathrm r}$};
\fill (8.2,0.5) circle (1.4pt); \node[right] at (8.2,0.5) {$A\mathbf x$};
\draw[line width=0.7pt, -{Latex[length=2mm]}] (1.7,0.55) .. controls (4.0,1.8) and (6.0,1.8) .. (8.1,0.55);
\node[above] at (4.9,1.5) {$A\mathbf x_{\mathrm r}=A\mathbf x$};
\fill (3.8,-1.2) circle (1.4pt); \node[left] at (3.75,-1.2) {$\mathbf x_{\mathrm n}$};
\draw[line width=0.7pt, -{Latex[length=2mm]}] (3.9,-1.15) .. controls (4.5,0.3) and (6.7,1.1) .. (7.34,0.1);
\node at (4.9,-1.2) {$A\mathbf x_{\mathrm n}=\mathbf 0$};
\node at (1.8,-1.85) {$\mathbb F^n$};
\node at (8.0,-1.85) {$\mathbb F^m$};
\end{tikzpicture}
\caption{The four fundamental subspaces of a matrix $A\in\mathbb F^{m\times n}$ of rank $r$. In the space $\mathbb F^n$ of the unknowns lie the row space and the null space, of dimensions $r$ and $n-r$; in the space $\mathbb F^m$ of the right-hand sides lie the column space and the left null space, of dimensions $r$ and $m-r$. For a real matrix, $\mathbf x=\mathbf x_{\mathrm r}+\mathbf x_{\mathrm n}$ splits into a part in the row space, which $A$ maps to $A\mathbf x$, and a part in the null space, which $A$ maps to $\mathbf 0$ (\cref{prop:row-null-real}). The two subspaces on each side meet only in $\mathbf 0$ and are drawn at right angles; Chapter~18 shows that each is the orthogonal complement of the other.}
\label{fig:four-subspaces}
\end{figure}

\section{Rank Revisited}
\label{sec:rank-revisited}

\subsection{The Many Faces of the Rank}

\Cref{ch:linear-systems} defined the rank by counting pivots, a definition that is convenient for computing by hand but tied to a particular algorithm. The following theorem collects descriptions that do not refer to elimination at all.

\begin{theorem}[Characterisations of the rank]
\label{thm:rank-characterisations}
For $A\in\mathbb F^{m\times n}$ the following numbers are equal:
\begin{enumerate}
\item the rank of $A$, the number of pivots of its reduced row echelon form;
\item the dimension of the column space, which is the largest number of linearly independent columns of $A$;
\item the dimension of the row space, which is the largest number of linearly independent rows of $A$;
\item the largest $k$ such that $A$ has an invertible $k\times k$ submatrix, formed by $k$ of its rows and $k$ of its columns ($k=0$ if $A=0$).
\end{enumerate}
\end{theorem}

\begin{proof}
Let $r\coloneqq\operatorname{rank}A$. By \cref{thm:four-dimensions}, $\dim\mathcal C(A)=\dim\mathcal C(A^{\mathsf T})=r$. Every independent set of columns lies in $\mathcal C(A)$ and has at most $\dim\mathcal C(A)$ elements, by \cref{thm:basis-extension}\,(i), and the pivot columns are $r$ independent columns; so the largest number of independent columns is $r$. The same argument applied to $A^{\mathsf T}$ gives the statement about rows.

(iv) Let $A_P\in\mathbb F^{m\times r}$ consist of the $r$ pivot columns of $A$. Its columns are independent, so $\operatorname{rank}A_P=r$ by \cref{prop:independence-columns}\,(i), and by what has just been shown, $A_P$ has $r$ independent rows. They form an $r\times r$ submatrix $S$ of $A$ with independent rows, so $\operatorname{rank}S=r$, and $S$ is invertible by \cref{thm:invertible}. Conversely, suppose that the submatrix $S$ in the rows $i_1,\ldots,i_k$ and the columns $j_1,\ldots,j_k$ is invertible. If $\sum_lx_l\mathbf a_{j_l}=\mathbf 0$, then in particular the components $i_1,\ldots,i_k$ of this vector vanish, which says $S\mathbf x=\mathbf 0$, so that $\mathbf x=\mathbf 0$. Hence the columns $j_1,\ldots,j_k$ of $A$ are independent, and $k\le r$.
\end{proof}

Description (ii) is the reason for the term \emph{column rank}, and (iii) for \emph{row rank}; the theorem says that the two coincide, which is \cref{thm:rank-transpose} once more. Description (iv) is the one that connects the rank with determinants, since a square matrix is invertible exactly when its determinant is nonzero (Chapter~17): the rank is the order of the largest nonvanishing minor. For the matrix $A$ of \cref{sec:four-subspaces}, the submatrix in the rows $1,2$ and the columns $1,3$, with the rows $(1,0)$ and $(2,1)$, is invertible, while all four $3\times3$ submatrices are singular, since $A$ has rank $2$.

The following special case is the fact about rank that the proof of \cref{thm:lagrange} used: if the rows of the derivative of the constraints are independent, some of its columns form an invertible matrix, to which the implicit function theorem applies.

\begin{corollary}[Matrices with independent rows]
\label{cor:independent-rows}
If the $m$ rows of $A\in\mathbb F^{m\times n}$ are linearly independent, then $\operatorname{rank}A=m\le n$, and some $m$ columns of $A$ form an invertible $m\times m$ matrix.
\end{corollary}

\begin{proof}
By \cref{thm:rank-characterisations}\,(iii), $\operatorname{rank}A=m$, and $m\le n$ by \cref{prop:rank-basic}\,(i). By (iv) there is an invertible $m\times m$ submatrix, which necessarily uses all $m$ rows; it consists of $m$ columns of $A$.
\end{proof}

\subsection{Rank Factorisation}

The bases of \cref{thm:four-bases} combine into a factorisation of $A$ that displays its rank. Let $C\in\mathbb F^{m\times r}$ be the matrix of the pivot columns of $A$, and let $R_0\in\mathbb F^{r\times n}$ consist of the $r$ nonzero rows of the reduced row echelon form $R$.

\begin{theorem}[Rank factorisation]
\label{thm:rank-factorisation}
Every $A\in\mathbb F^{m\times n}$ of rank $r$ satisfies
\[
A=CR_0,
\]
where $C\in\mathbb F^{m\times r}$ has independent columns, a basis of $\mathcal C(A)$, and $R_0\in\mathbb F^{r\times n}$ has independent rows, a basis of $\mathcal C(A^{\mathsf T})$. Consequently $A$ is a sum of $r$ matrices of rank one. Conversely, if $A=BD$ with $B\in\mathbb F^{m\times k}$ and $D\in\mathbb F^{k\times n}$, then $k\ge r$, and $A$ is not a sum of fewer than $r$ matrices of rank one.
\end{theorem}

\begin{proof}
By the column reading of the product (\cref{thm:four-readings}\,(ii)), the $j$th column of $CR_0$ is the combination of the columns $\mathbf a_{p_1},\ldots,\mathbf a_{p_r}$ of $C$ with the entries of the $j$th column of $R_0$ as coefficients. For a pivot column $j=p_k$, this column of $R_0$ is $\mathbf e_k\in\mathbb F^r$, and the combination is $\mathbf a_{p_k}$. For a free column $j=f_k$, the coefficients are $r_{1f_k},\ldots,r_{rf_k}$, and the combination is $\mathbf a_{f_k}$ by \cref{thm:four-bases}\,(i). So $CR_0=A$. The properties of $C$ and $R_0$ are (i) and (ii) of \cref{thm:four-bases}. By the outer product reading (\cref{thm:four-readings}\,(iv)), $CR_0=\sum_{i=1}^r\mathbf c_i\boldsymbol\rho_i^{\mathsf T}$, where $\mathbf c_i$ is the $i$th column of $C$ and $\boldsymbol\rho_i^{\mathsf T}$ the $i$th row of $R_0$, and each term is a nonzero outer product, which has rank one (\cref{sec:rank}). If $A=BD$, then every column of $A$ is a combination of the $k$ columns of $B$, by the column reading, so $\mathcal C(A)\subseteq\mathcal C(B)$ and $r=\dim\mathcal C(A)\le\dim\mathcal C(B)\le k$. A sum $\sum_{i=1}^k\mathbf u_i\mathbf v_i^{\mathsf T}$ of $k$ outer products is such a product, with the columns $\mathbf u_i$ and the rows $\mathbf v_i^{\mathsf T}$.
\end{proof}

For the matrix of \cref{sec:four-subspaces}, the factorisation and the sum of rank-one matrices are
\[
\begin{bmatrix}1&2&0&3\\2&4&1&8\\-1&-2&1&-1\end{bmatrix}
=\begin{bmatrix}1&0\\2&1\\-1&1\end{bmatrix}\begin{bmatrix}1&2&0&3\\0&0&1&2\end{bmatrix}
=\begin{bmatrix}1\\2\\-1\end{bmatrix}\begin{bmatrix}1&2&0&3\end{bmatrix}+\begin{bmatrix}0\\1\\1\end{bmatrix}\begin{bmatrix}0&0&1&2\end{bmatrix}.
\]
The factorisation stores the $12$ entries of $A$ in $6+8=14$ numbers, which saves nothing for so small a matrix; but an $m\times n$ matrix of rank $r$ needs only $r(m+n)$ numbers in factored form instead of $mn$, and a product with a vector costs about $2r(m+n)$ operations instead of $2mn$. For $m=n=1000$ and $r=10$ this is a saving by a factor of $50$. Matrices of low rank, or close to low rank, are common in applications, and Chapter~21 shows how the singular value decomposition finds the best approximation of a given rank. Strang and Moler propose the factorisation $A=CR_0$, which they call $A=CR$, as the starting point of a first course~\cite{strang2022}.

For $r=1$ the theorem says that \emph{a matrix has rank one if and only if it is a nonzero outer product $\mathbf u\mathbf v^{\mathsf T}$}: one direction is the factorisation with $C=\mathbf u$ and $R_0=\mathbf v^{\mathsf T}$, and the other is the example of \cref{sec:rank}. All columns of a rank-one matrix are multiples of $\mathbf u$, and all rows are multiples of $\mathbf v^{\mathsf T}$.

\subsection{Sums and Products}

\begin{proposition}[Rank of sums and products]
\label{prop:rank-sum}
Let $A,B\in\mathbb F^{m\times n}$ and $D\in\mathbb F^{n\times p}$.
\begin{enumerate}
\item $\mathcal C(A+B)\subseteq\mathcal C(A)+\mathcal C(B)$, and $\operatorname{rank}(A+B)\le\operatorname{rank}A+\operatorname{rank}B$.
\item $|\operatorname{rank}A-\operatorname{rank}B|\le\operatorname{rank}(A-B)$.
\item $\mathcal C(AD)\subseteq\mathcal C(A)$, and $\operatorname{rank}(AD)\le\operatorname{rank}A$.
\end{enumerate}
\end{proposition}

\begin{proof}
(i) $(A+B)\mathbf x=A\mathbf x+B\mathbf x$ lies in $\mathcal C(A)+\mathcal C(B)$, and by \cref{thm:dimension-formula} this sum has dimension at most $\dim\mathcal C(A)+\dim\mathcal C(B)$. By \cref{prop:subspace-dimension} and \cref{thm:rank-characterisations}, $\operatorname{rank}(A+B)=\dim\mathcal C(A+B)\le\operatorname{rank}A+\operatorname{rank}B$. (ii) By (i), $\operatorname{rank}A=\operatorname{rank}((A-B)+B)\le\operatorname{rank}(A-B)+\operatorname{rank}B$, and in the same way $\operatorname{rank}B\le\operatorname{rank}(B-A)+\operatorname{rank}A$, where $\operatorname{rank}(B-A)=\operatorname{rank}(A-B)$ since the two matrices have the same column space. (iii) The columns of $AD$ are combinations of the columns of $A$, by the column reading of the product, which \cref{ch:vectors-matrices} anticipated would serve to compare ranks; so $\mathcal C(AD)\subseteq\mathcal C(A)$, and the dimensions compare accordingly.
\end{proof}

By (ii), adding a matrix of rank one, as in the Sherman--Morrison formula (\cref{thm:sherman-morrison}) or in a step of Gaussian elimination, changes the rank by at most one. Part (iii) is the first inequality of \cref{prop:rank-product}, now with a proof by dimensions instead of elimination; the second inequality, $\operatorname{rank}(AD)\le\operatorname{rank}D$, which \cref{ch:linear-systems} derived from the equality of the ranks of a matrix and its transpose, is now fully proved, since that equality is. Chapter~16 proves it again from the null spaces, and sharpens both inequalities to Sylvester's inequality $\operatorname{rank}A+\operatorname{rank}D-n\le\operatorname{rank}(AD)$.

\subsection{Full Rank and One-Sided Inverses}

A matrix $A\in\mathbb F^{m\times n}$ has rank at most $\min(m,n)$, and it has \term{full column rank} if $\operatorname{rank}A=n$ and \term{full row rank} if $\operatorname{rank}A=m$. These two conditions are the abstract forms of the questions of uniqueness and existence for linear systems, as the following theorem shows; recall that the map $\mathbf x\mapsto A\mathbf x$ is \term{injective} if different vectors have different images, and \term{surjective} if every vector of $\mathbb F^m$ is an image (\cref{sec:injections}).

\begin{theorem}[Full rank]
\label{thm:full-rank}
Let $A\in\mathbb F^{m\times n}$.
\begin{enumerate}
\item The following are equivalent: $\operatorname{rank}A=n$; the columns of $A$ are linearly independent; $\mathcal N(A)=\{\mathbf 0\}$; the map $\mathbf x\mapsto A\mathbf x$ is injective; $A\mathbf x=\mathbf b$ has at most one solution for every $\mathbf b$; $A$ has a left inverse, $BA=I_n$ for some $B\in\mathbb F^{n\times m}$.
\item The following are equivalent: $\operatorname{rank}A=m$; the columns of $A$ span $\mathbb F^m$; $\mathcal N(A^{\mathsf T})=\{\mathbf 0\}$; the map $\mathbf x\mapsto A\mathbf x$ is surjective; $A\mathbf x=\mathbf b$ has at least one solution for every $\mathbf b$; $A$ has a right inverse, $AC=I_m$ for some $C\in\mathbb F^{n\times m}$.
\end{enumerate}
\end{theorem}

\begin{proof}
(i) The first three conditions are equivalent by \cref{prop:independence-columns}\,(i). The map is injective if and only if $\mathcal N(A)=\{\mathbf 0\}$, since $A\mathbf x=A\mathbf x'$ means $A(\mathbf x-\mathbf x')=\mathbf 0$, and the fifth condition says the same in other words. A left inverse makes the map injective by \cref{prop:one-sided-inverses}\,(ii), whose proof holds for complex matrices as well. Conversely, let $\operatorname{rank}A=n$. Then $A^{\mathsf T}\in\mathbb F^{n\times m}$ has rank $n$ by \cref{thm:four-dimensions}, which is its number of rows, so by \cref{thm:rank-solutions}\,(iii) the systems $A^{\mathsf T}\mathbf z_i=\mathbf e_i$, $i=1,\ldots,n$, have solutions $\mathbf z_i\in\mathbb F^m$. The matrix $Z$ with the columns $\mathbf z_i$ satisfies $A^{\mathsf T}Z=I_n$, and $B\coloneqq Z^{\mathsf T}$ satisfies $BA=(A^{\mathsf T}Z)^{\mathsf T}=I_n$.

(ii) The first two conditions are equivalent by \cref{prop:independence-columns}\,(ii), and the first and the third by $\dim\mathcal N(A^{\mathsf T})=m-\operatorname{rank}A$ (\cref{thm:four-dimensions}). The second, fourth and fifth conditions all say that $\mathcal C(A)=\mathbb F^m$. A right inverse $C$ gives the solution $\mathbf x=C\mathbf b$ for every $\mathbf b$ (\cref{prop:one-sided-inverses}\,(iii)). Conversely, if every system is solvable, the solutions $\mathbf c_k$ of $A\mathbf c_k=\mathbf e_k$ for $k=1,\ldots,m$ form the columns of a right inverse $C$.
\end{proof}

For a square matrix, $m=n$, the two lists merge, and with \cref{thm:invertible} we obtain a result that \cref{ch:vectors-matrices} anticipated after \cref{thm:one-sided-square}: \emph{a square matrix is injective if and only if it is surjective}. Indeed, by the rank--nullity theorem $\dim\mathcal C(A)=n-\dim\mathcal N(A)$, so $\mathcal N(A)=\{\mathbf 0\}$ exactly when $\mathcal C(A)$ has dimension $n$, that is, when $\mathcal C(A)=\mathbb F^n$ (\cref{prop:subspace-dimension}). The argument counts dimensions, and it is exactly here that the shift operators of \cref{sec:inverse} escape: the right shift on the space of sequences is injective but not surjective, and there is no finite dimension to count. Chapter~16 proves the same statement for linear maps of any finite-dimensional space to itself.

A tall matrix of full column rank has many left inverses, and a wide matrix of full row rank many right inverses. For the column $\mathbf a=(1,2)$, the left inverses are the rows $(b_1,b_2)$ with $b_1+2b_2=1$, among them $(1,0)$, $(0,1/2)$ and $(1,2)/5$, the last of which \cref{sec:inverse} described as the left inverse given by orthogonal projection. For real matrices there is a natural choice, given by the following proposition; the example $(1,2)/5=(\mathbf a^{\mathsf T}\mathbf a)^{-1}\mathbf a^{\mathsf T}$ is its simplest case.

\begin{proposition}[The Gram matrix and full column rank]
\label{prop:gram-rank}
For $A\in\mathbb R^{m\times n}$, $\mathcal N(A^{\mathsf T}A)=\mathcal N(A)$ and $\operatorname{rank}(A^{\mathsf T}A)=\operatorname{rank}A$. In particular, $A^{\mathsf T}A$ is invertible if and only if $A$ has full column rank, and then $(A^{\mathsf T}A)^{-1}A^{\mathsf T}$ is a left inverse of $A$. Similarly, if $A$ has full row rank, $AA^{\mathsf T}$ is invertible, and $A^{\mathsf T}(AA^{\mathsf T})^{-1}$ is a right inverse of $A$.
\end{proposition}

\begin{proof}
If $A\mathbf x=\mathbf 0$, then $A^{\mathsf T}A\mathbf x=\mathbf 0$. Conversely, if $A^{\mathsf T}A\mathbf x=\mathbf 0$, then $\|A\mathbf x\|^2=\langle\mathbf x,A^{\mathsf T}A\mathbf x\rangle=0$ by \cref{prop:gram}, so $A\mathbf x=\mathbf 0$. The two matrices have $n$ columns and the same null space, so they have the same rank by the rank--nullity theorem. The square matrix $A^{\mathsf T}A$ is invertible if and only if its rank is $n$ (\cref{thm:invertible}), that is, if $\operatorname{rank}A=n$, and then $(A^{\mathsf T}A)^{-1}A^{\mathsf T}A=I$. The second statement follows by applying the first to $A^{\mathsf T}$, which has full column rank when $A$ has full row rank, and transposing.
\end{proof}

The left inverse $(A^{\mathsf T}A)^{-1}A^{\mathsf T}$ is the matrix of the normal equations of least squares (\cref{sec:extrema-n}): if $A$ has full column rank, the system $A^{\mathsf T}A\mathbf x=A^{\mathsf T}\mathbf b$ has exactly one solution for every $\mathbf b$, whether or not $A\mathbf x=\mathbf b$ is consistent. Chapter~18 shows that this solution minimises $\|A\mathbf x-\mathbf b\|$, and Chapter~21 shows that both formulas of the proposition are special cases of the pseudoinverse. For complex matrices the proposition holds with $A^{\mathsf H}$ in place of $A^{\mathsf T}$, for the reason explained after \cref{prop:row-null-real}.

\section{Solvability and the Fredholm Alternative}
\label{sec:fredholm}

\subsection{The Alternative}

The system $A\mathbf x=\mathbf b$ is consistent if and only if $\mathbf b\in\mathcal C(A)$. This is a reformulation, as \cref{sec:systems} remarked, and the question is how to recognise the vectors of the column space. The left null space answers it: the vectors of the column space are those that satisfy the linear relations recorded by the left null space.

\begin{theorem}[Fredholm alternative]
\label{thm:fredholm}
Let $A\in\mathbb F^{m\times n}$ and $\mathbf b\in\mathbb F^m$. Then $A\mathbf x=\mathbf b$ has a solution if and only if
\[
\mathbf y^{\mathsf T}\mathbf b=0\qquad\text{for every }\mathbf y\in\mathcal N(A^{\mathsf T}).
\]
Equivalently, exactly one of the following two statements holds:
\begin{enumerate}
\item the system $A\mathbf x=\mathbf b$ has a solution;
\item there is a vector $\mathbf y\in\mathbb F^m$ with $\mathbf y^{\mathsf T}A=\mathbf 0^{\mathsf T}$ and $\mathbf y^{\mathsf T}\mathbf b=1$.
\end{enumerate}
\end{theorem}

\begin{proof}
If $A\mathbf x=\mathbf b$ and $\mathbf y^{\mathsf T}A=\mathbf 0^{\mathsf T}$, then $\mathbf y^{\mathsf T}\mathbf b=\mathbf y^{\mathsf T}A\mathbf x=0$. For the converse, let $r=\operatorname{rank}A$, and let $Y\in\mathbb F^{m\times(m-r)}$ be a matrix whose columns form a basis of $\mathcal N(A^{\mathsf T})$, which has dimension $m-r$ by \cref{thm:four-dimensions}. The vectors $\mathbf b$ with $\mathbf y^{\mathsf T}\mathbf b=0$ for all $\mathbf y\in\mathcal N(A^{\mathsf T})$ are those with $Y^{\mathsf T}\mathbf b=\mathbf 0$, since every $\mathbf y$ in the left null space is a combination of the columns of $Y$; so they form the subspace $\mathcal N(Y^{\mathsf T})$. By the first part, $\mathcal C(A)\subseteq\mathcal N(Y^{\mathsf T})$. The matrix $Y^{\mathsf T}\in\mathbb F^{(m-r)\times m}$ has rank $m-r$, since its rows, the columns of $Y$, are independent (\cref{thm:rank-characterisations}), and so $\dim\mathcal N(Y^{\mathsf T})=m-(m-r)=r=\dim\mathcal C(A)$ by \cref{thm:four-dimensions}. A subspace of the same dimension is the whole space (\cref{prop:subspace-dimension}), so $\mathcal C(A)=\mathcal N(Y^{\mathsf T})$, which is the claim.

For the alternative: if (i) holds, then $\mathbf y^{\mathsf T}\mathbf b=0$ for every $\mathbf y$ in the left null space, so (ii) fails. If (i) fails, there is $\mathbf y\in\mathcal N(A^{\mathsf T})$ with $c\coloneqq\mathbf y^{\mathsf T}\mathbf b\ne0$, and $\mathbf y/c$ satisfies (ii).
\end{proof}

The proof shows that it suffices to test $\mathbf b$ against a \emph{basis} of the left null space, that is, $m-r$ conditions, one for each dependent row of $A$. Elimination produces such a basis, the last $m-r$ rows of $G$ (\cref{thm:four-bases}\,(iv)), and the test is then exactly the test of \cref{thm:solvability}: the components $d_{r+1},\ldots,d_m$ of $\mathbf d=G\mathbf b$, which must vanish for a consistent system, are the products of $\mathbf b$ with the last rows of $G$. For the matrix of \cref{sec:four-subspaces}, with the left null space spanned by $(3,-1,1)$, the system is consistent if and only if $3b_1-b_2+b_3=0$. In words, \emph{a system is consistent if and only if every linear relation among its equations also holds among the right-hand sides}: if a combination of the rows of $A$ vanishes, the same combination of the components of $\mathbf b$ must vanish, and this necessary condition is also sufficient. A vector $\mathbf y$ as in (ii) is a \emph{certificate} of inconsistency: whoever claims that a system has no solution can prove it by exhibiting $\mathbf y$, and anyone can check the proof by computing $\mathbf y^{\mathsf T}A$ and $\mathbf y^{\mathsf T}\mathbf b$, without repeating the elimination.

For a square matrix the theorem takes the form in which Fredholm stated it for integral equations~\cite{fredholm1903}.

\begin{corollary}[Fredholm alternative for square matrices]
\label{cor:fredholm-square}
For $A\in\mathbb F^{n\times n}$, exactly one of the following holds:
\begin{enumerate}
\item $A\mathbf x=\mathbf b$ has exactly one solution for every $\mathbf b\in\mathbb F^n$;
\item the homogeneous systems $A\mathbf x=\mathbf 0$ and $A^{\mathsf T}\mathbf y=\mathbf 0$ have nonzero solutions, the spaces of their solutions have the same dimension, and $A\mathbf x=\mathbf b$ has a solution if and only if $\mathbf y^{\mathsf T}\mathbf b=0$ for all solutions $\mathbf y$ of $A^{\mathsf T}\mathbf y=\mathbf 0$.
\end{enumerate}
\end{corollary}

\begin{proof}
If $\operatorname{rank}A=n$, then (i) holds by \cref{thm:invertible}, and $\mathcal N(A)=\{\mathbf 0\}$, so (ii) fails. If $\operatorname{rank}A=r<n$, then (i) fails, and $\dim\mathcal N(A)=\dim\mathcal N(A^{\mathsf T})=n-r>0$ by \cref{thm:four-dimensions}; the last statement of (ii) is \cref{thm:fredholm}.
\end{proof}

The corollary generalises \cref{prop:singular-image}, which found for a singular square matrix one vector $\mathbf y\ne\mathbf 0$ with $\mathbf y^{\mathsf T}A=\mathbf 0^{\mathsf T}$ and concluded that the image lies in a hyperplane. Now we know the image exactly: it is the set of vectors $\mathbf b$ with $\mathbf y^{\mathsf T}\mathbf b=0$ for all such $\mathbf y$. For real matrices the condition says that $\mathbf b$ is orthogonal to the left null space, and \cref{prop:row-null-real}\,(ii) is the same statement in the form of a direct sum; Chapter~18 phrases it as $\mathcal C(A)=\mathcal N(A^{\mathsf T})^\perp$. The same structure, with infinite-dimensional spaces of functions in place of $\mathbb F^n$, governs integral equations and boundary value problems: an equation such as $-u''=f$ on $[0,\pi]$ with $u'(0)=u'(\pi)=0$ has a solution only if $\int_0^\pi f=0$, because the constant functions solve the homogeneous problem, and the condition is the Fredholm condition. Volume~II returns to this in the setting of Hilbert spaces.

\subsection{The Incidence Matrix of a Graph}

The four subspaces have a vivid meaning for the matrices that describe graphs. Let $G=(V,E)$ be a graph in the sense of \cref{def:graph}, with the vertices $1,\ldots,n$ and the edges $e_1,\ldots,e_m$, and give every edge a direction, from one of its endpoints, its tail, to the other, its head; which direction is chosen does not matter. The \term{incidence matrix} of the directed graph is the matrix $A\in\mathbb R^{m\times n}$ with one row for each edge and one column for each vertex, in which the row of the edge $e_k$ from $u$ to $v$ has the entry $-1$ in column $u$, the entry $+1$ in column $v$ and zeros elsewhere. The graph of \cref{fig:incidence} has $n=4$ vertices and $m=5$ edges, and its incidence matrix, with the edges listed beside its rows, is
\[
A=\begin{bmatrix}-1&1&0&0\\-1&0&1&0\\0&-1&1&0\\0&-1&0&1\\0&0&-1&1\end{bmatrix}
\begin{matrix}e_1\colon1\to2\\e_2\colon1\to3\\e_3\colon2\to3\\e_4\colon2\to4\\e_5\colon3\to4\end{matrix}
\]

\begin{figure}
\sdfloat{\begin{tikzpicture}[font=\scriptsize, x=1.1cm, y=1.1cm,
  gv/.style={circle, draw, line width=0.5pt, fill=white, inner sep=1.2pt, minimum size=13pt},
  ga/.style={line width=0.7pt, -{Latex[length=2.2mm]}, shorten >=1pt}]
\node[gv] (1) at (0,1) {$1$};
\node[gv] (2) at (1.7,2) {$2$};
\node[gv] (3) at (1.7,0) {$3$};
\node[gv] (4) at (3.4,1) {$4$};
\draw[ga] (1) -- node[above left=-1pt] {$e_1$} (2);
\draw[ga] (1) -- node[below left=-1pt] {$e_2$} (3);
\draw[ga] (2) -- node[left] {$e_3$} (3);
\draw[ga] (2) -- node[above right=-1pt] {$e_4$} (4);
\draw[ga] (3) -- node[below right=-1pt] {$e_5$} (4);
\end{tikzpicture}}{A directed graph with $4$ vertices and $5$ edges. Its incidence matrix has rank $3=n-1$, since the graph is connected; its null space consists of the constant vectors, and its left null space has dimension $m-n+1=2$, spanned by the two triangles $e_1,e_3,e_2$ and $e_3,e_5,e_4$ traversed as loops.\label{fig:incidence}}
\end{figure}

The product $A\mathbf x$ has a direct interpretation. Think of $x_v$ as a potential at the vertex $v$, for instance the voltage at a node of an electrical network or the height of a point in a landscape. Then the $k$th component of $A\mathbf x$ is $x_v-x_u$ for the edge $e_k$ from $u$ to $v$, the potential difference along the edge. Similarly, think of $y_k$ as a flow along the edge $e_k$, positive in the direction of the edge, such as an electric current. The component of $A^{\mathsf T}\mathbf y$ for the vertex $v$ is the sum of the $y_k$ over the edges that end at $v$ minus the sum over the edges that start at $v$, the net flow into $v$. The four subspaces now have the following meaning. We write $c$ for the number of connected components of the graph (\cref{sec:graphs}).

\begin{proposition}[The four subspaces of an incidence matrix]
\label{prop:incidence}
Let $A\in\mathbb R^{m\times n}$ be the incidence matrix of a graph with $n$ vertices, $m$ edges and $c$ connected components, with its edges directed arbitrarily.
\begin{enumerate}
\item $\mathcal N(A)$ consists of the vectors that are constant on each component, and $\dim\mathcal N(A)=c$.
\item $\operatorname{rank}A=n-c$.
\item $\dim\mathcal N(A^{\mathsf T})=m-n+c$. The vectors of $\mathcal N(A^{\mathsf T})$ are the flows that are \term{conserved} at every vertex: the net flow into each vertex is zero.
\item A vector $\mathbf b\in\mathbb R^m$ lies in $\mathcal C(A)$, that is, it consists of the potential differences of some potential, if and only if $\mathbf y^{\mathsf T}\mathbf b=0$ for every conserved flow $\mathbf y$.
\end{enumerate}
\end{proposition}

\begin{proof}
(i) $A\mathbf x=\mathbf 0$ means $x_u=x_v$ for every edge $uv$. If $u$ and $w$ lie in the same component, there is a walk $u=v_0,v_1,\ldots,v_k=w$ (\cref{def:walk}), and $x$ takes the same value at consecutive vertices of the walk, hence $x_u=x_w$. Conversely, a vector that is constant on each component satisfies $x_u=x_v$ for every edge, since the endpoints of an edge lie in the same component. Let $\mathbf 1_{K_1},\ldots,\mathbf 1_{K_c}\in\mathbb R^n$ be the indicator vectors of the components $K_1,\ldots,K_c$, with the entry $1$ at the vertices of $K_j$ and $0$ elsewhere. They lie in $\mathcal N(A)$, they span it, since a vector constant on each component is $\sum_jx_{v_j}\mathbf 1_{K_j}$ for any choice of vertices $v_j\in K_j$, and they are independent, since the components are disjoint and nonempty, so that the coefficient of $\mathbf 1_{K_j}$ in a combination is its entry at a vertex of $K_j$. So $\dim\mathcal N(A)=c$. (ii) and (iii) follow from \cref{thm:four-dimensions}: $\operatorname{rank}A=n-\dim\mathcal N(A)=n-c$ and $\dim\mathcal N(A^{\mathsf T})=m-(n-c)$. The description of $\mathcal N(A^{\mathsf T})$ is the interpretation of $A^{\mathsf T}\mathbf y$ given above. (iv) is \cref{thm:fredholm}.
\end{proof}

Part (iii) is \term{Kirchhoff's current law}, which says that the currents in an electrical network are conserved at every node, and the number $m-n+c$ is the number of independent \emph{loops} of the network, as the next proposition explains. Part (iv) contains \term{Kirchhoff's voltage law}: the voltages along the edges come from potentials at the nodes if and only if they add up to zero around every loop. Kirchhoff's paper of 1847, which computes the currents of a network from these two laws, counts the independent loops in this way and introduces what are now called spanning trees~\cite{kirchhoff1847}.

A cycle of the graph (\cref{sec:graphs}) gives a conserved flow. Let $(v_0,v_1,\ldots,v_{k-1},v_0)$ be a cycle, traversed in this order, and define $\mathbf y\in\mathbb R^m$ by $y_e\coloneqq+1$ for an edge $e$ of the cycle whose direction agrees with the traversal, $y_e\coloneqq-1$ for an edge of the cycle whose direction is opposite, and $y_e\coloneqq0$ for the edges not on the cycle. This is the \term{cycle vector} of the traversed cycle, a unit flow around the cycle. At each vertex $v_i$ of the cycle, the flow arrives along the cycle edge from $v_{i-1}$ and leaves along the cycle edge to $v_{i+1}$ (indices modulo $k$), and whatever the directions of these two edges, the first contributes $+1$ and the second $-1$ to the net flow into $v_i$; at the other vertices nothing flows. So every cycle vector lies in $\mathcal N(A^{\mathsf T})$. In the example of \cref{fig:incidence}, the triangle $1\to2\to3\to1$ uses $e_1$ and $e_3$ forwards and $e_2$ backwards, and the triangle $2\to3\to4\to2$ uses $e_3$ and $e_5$ forwards and $e_4$ backwards, which gives the cycle vectors
\[
\mathbf y_1=(1,-1,1,0,0),\qquad\mathbf y_2=(0,0,1,-1,1),
\]
and one checks that $\mathbf y_1^{\mathsf T}A=\mathbf y_2^{\mathsf T}A=\mathbf 0^{\mathsf T}$. Neither is a multiple of the other, so they form a basis of $\mathcal N(A^{\mathsf T})$, whose dimension is $5-4+1=2$. The outer cycle $1\to2\to4\to3\to1$ has the cycle vector $(1,-1,0,1,-1)=\mathbf y_1-\mathbf y_2$, a combination of the two. By (iv), a vector $\mathbf b$ of edge voltages comes from node potentials if and only if $b_1-b_2+b_3=0$ and $b_3-b_4+b_5=0$.

\begin{proposition}[A basis of cycles]
\label{prop:cycle-basis}
Every graph has a basis of $\mathcal N(A^{\mathsf T})$ consisting of $m-n+c$ cycle vectors. In particular, the cycle vectors span the conserved flows, and $\mathcal N(A^{\mathsf T})=\{\mathbf 0\}$ if and only if the graph is a forest.
\end{proposition}

\begin{proof}
Each component $K_j$, with $n_j$ vertices, is connected and has a spanning tree with $n_j-1$ edges, by \cref{cor:spanning-tree,thm:tree}; together these trees form a spanning forest $F$ with $\sum_j(n_j-1)=n-c$ edges. Let $e$ be one of the $m-(n-c)$ edges not in $F$, with endpoints $u$ and $v$. They lie in the same component, and by \cref{thm:tree}\,(iv) the tree of this component contains exactly one path from $v$ to $u$; it has length at least $2$, since $e$ itself is not in the tree. The path followed by $e$ is a cycle $C_e$, which we traverse so that $e$ is used in its own direction. Its cycle vector $\mathbf z_e$ lies in $\mathcal N(A^{\mathsf T})$; its entry for the edge $e$ is $+1$, and its entries for the other edges outside $F$ are $0$, since $C_e$ consists of $e$ and edges of $F$. A combination $\sum_ea_e\mathbf z_e=\mathbf 0$ has the entry $a_e$ at the edge $e$, so all $a_e$ vanish, and the $m-n+c$ vectors $\mathbf z_e$ are independent. By \cref{prop:incidence}\,(iii) and \cref{thm:basis-extension}\,(iii), they form a basis. The space is $\{\mathbf 0\}$ if and only if $m=n-c$, which by \cref{lem:components-edges} characterises forests.
\end{proof}

The proof gives a second derivation of the count of \cref{lem:components-edges}: a graph with $n$ vertices, $m$ edges and $c$ components has $m\ge n-c$, because $m-n+c$ is a dimension. The product $A^{\mathsf T}A\in\mathbb R^{n\times n}$ is the \term{Laplacian matrix} of the graph, with the degrees of the vertices on its diagonal and the entry $-1$ for each pair of adjacent vertices; for the graph of \cref{fig:incidence},
\[
A^{\mathsf T}A=\begin{bmatrix}2&-1&-1&0\\-1&3&-1&-1\\-1&-1&3&-1\\0&-1&-1&2\end{bmatrix}.
\]
By \cref{prop:gram-rank}, its null space is that of $A$, the constant vectors on each component, so it has rank $n-c$. The Laplacian matrix governs electrical networks, random walks on graphs and the partitioning of graphs, and it is one of the main objects of Volume~IX; Biggs~\cite{biggs1993} and Godsil and Royle~\cite{godsil2001} develop the algebraic theory of graphs that begins here.

\section{Computing Bases and Rank: A First Look}
\label{sec:computing-rank}

\subsection{An Algorithm for the Four Subspaces}

\Cref{thm:four-bases} is an algorithm in disguise: one elimination produces bases of all four subspaces. To obtain the matrix $G$ of the left null space along with $R$, every row operation applied to $A$ is applied to the identity matrix as well, as in the example of \cref{sec:four-subspaces}; equivalently, the first $n$ columns of the augmented matrix $[A\mid I_m]$ are reduced by \cref{alg:rref}. \Cref{alg:four-subspaces} states the procedure.

\begin{algorithm}
\begin{algorithmic}[1]
\Require $A=(a_{ij})\in\mathbb F^{m\times n}$
\Ensure the rank $r$ and bases of $\mathcal C(A)$, $\mathcal C(A^{\mathsf T})$, $\mathcal N(A)$ and $\mathcal N(A^{\mathsf T})$
\Name{FourSubspaces}{A}
\State $[R\mid G]\gets[A\mid I_m]$
\State reduce the first $n$ columns of $[R\mid G]$ by \Call{ReducedEchelon}{}, applying each row operation to whole rows \Comment{now $GA=R$}
\State $r\gets$ the number of pivots; $p_1<\cdots<p_r$ the pivot columns, $f_1<\cdots<f_{n-r}$ the free columns
\State $\mathcal C(A)$: the columns $\mathbf a_{p_1},\ldots,\mathbf a_{p_r}$ of the \emph{original} $A$
\State $\mathcal C(A^{\mathsf T})$: the rows $1,\ldots,r$ of $R$, as columns
\For{$k=1,\ldots,n-r$}
  \State $\mathbf s_k\gets\mathbf e_{f_k}-\sum_{i=1}^rr_{if_k}\mathbf e_{p_i}$ \Comment{the special solution of \cref{thm:homogeneous-solutions}}
\EndFor
\State $\mathcal N(A)$: $\mathbf s_1,\ldots,\mathbf s_{n-r}$
\State $\mathcal N(A^{\mathsf T})$: the rows $r+1,\ldots,m$ of $G$, as columns
\State \Return $r$ and the four bases
\end{algorithmic}
\caption{Bases of the four fundamental subspaces by Gauss--Jordan elimination (\cref{thm:four-bases}). The reduction performs at most $r$ pivot steps, each of which updates $m-1$ rows of length $m+n$, so it costs at most about $2rm(m+n)$ operations.}
\label{alg:four-subspaces}
\end{algorithm}

In exact arithmetic, over the rational numbers or a finite field, the algorithm is all one needs, and computer algebra systems use it in this form. In floating-point arithmetic it is unreliable for a reason that has nothing to do with the algorithm and everything to do with the rank: the decision whether a pivot is zero cannot be made reliably when the entries carry rounding errors, and the following proposition shows that the answer is extremely sensitive to such errors. We state it for real matrices and the norm $\|\cdot\|_\infty$ of \cref{sec:trace-norms}.

\begin{proposition}[The rank under small perturbations]
\label{prop:rank-perturbation}
Let $A\in\mathbb R^{m\times n}$ have rank $r$.
\begin{enumerate}
\item There is a $\delta>0$ such that $\operatorname{rank}(A+E)\ge r$ for every $E\in\mathbb R^{m\times n}$ with $\|E\|_\infty<\delta$: small perturbations do not decrease the rank.
\item If $r<\min(m,n)$, then for every $\varepsilon>0$ there is an $E\in\mathbb R^{m\times n}$ with $\|E\|_\infty\le\varepsilon$ and $\operatorname{rank}(A+E)=\min(m,n)$: arbitrarily small perturbations can increase the rank to its largest possible value.
\end{enumerate}
\end{proposition}

\begin{proof}
(i) For $r=0$ there is nothing to prove. Otherwise, by \cref{thm:rank-characterisations}\,(iv), $A$ has an invertible $r\times r$ submatrix $S$. Let $E_S$ be the submatrix of $E$ in the same rows and columns. Each row sum of absolute values of $E_S$ is part of a row sum of $E$, so $\|E_S\|_\infty\le\|E\|_\infty$ by \cref{thm:norm-formulas}. With $\delta\coloneqq1/\|S^{-1}\|_\infty$, \cref{cor:perturbed-inverse} shows that $S+E_S$ is invertible whenever $\|E\|_\infty<\delta$, and $S+E_S$ is a submatrix of $A+E$, which therefore has rank at least $r$.

(ii) Let $B$ be a matrix with an invertible $k\times k$ submatrix $S$ in the rows $I$ and the columns $J$, where $k<\min(m,n)$; initially $B=A$ and $k=r$. Choose a row $i\notin I$ and a column $j\notin J$, and consider the $(k+1)\times(k+1)$ submatrix of $B+tE_{ij}$ in the rows $I\cup\{i\}$ and the columns $J\cup\{j\}$. Ordered so that $S$ comes first, it has the block form
\[
T(t)=\begin{bmatrix}S&\mathbf u\\\mathbf w^{\mathsf T}&d+t\end{bmatrix}
\]
with $\mathbf u$, $\mathbf w$ and $d=b_{ij}$ taken from $B$, and by \cref{thm:schur-complement}\,(ii), $T(t)$ is invertible if and only if its $1\times1$ Schur complement $d+t-\mathbf w^{\mathsf T}S^{-1}\mathbf u$ is nonzero. This fails for at most one value of $t$, so we can choose $t$ with $0<|t|\le\varepsilon$ and $T(t)$ invertible. Then $B+tE_{ij}$ has an invertible submatrix of order $k+1$. Repeating the step $\min(m,n)-r$ times gives a matrix $A+E$ with an invertible submatrix of order $\min(m,n)$, hence of rank $\min(m,n)$, by \cref{thm:rank-characterisations}\,(iv) and \cref{prop:rank-basic}\,(i). Each step changes a single entry, by at most $\varepsilon$ in absolute value, in a row that is not changed again, since it joins $I$; so every row of $E$ has at most one nonzero entry, and $\|E\|_\infty\le\varepsilon$.
\end{proof}

So the matrices of full rank form an \emph{open} set: a matrix of full rank keeps it under all sufficiently small perturbations. The matrices of lower rank, on the other hand, are exceptional: every one of them is the limit of matrices of full rank. In the language of analysis, the rank is a lower semicontinuous function of the entries, and it is continuous exactly at the matrices of full rank. The simplest example is $\operatorname{diag}(1,\varepsilon)$, which has rank $2$ for every $\varepsilon\ne0$ and rank $1$ for $\varepsilon=0$. For computation the consequence is serious. Rounding the entries of a matrix to floating-point numbers, or rounding the results of the operations of elimination, is a small perturbation, and it will typically turn a matrix of low rank into one of full rank; a pivot that should be zero then comes out as a tiny nonzero number. Conversely, a pivot that should be tiny but nonzero may come out as exactly zero.

\subsection{An Experiment}

\Cref{tab:rank-experiment} compares three ways of computing the rank of some matrices stored in binary64 floating-point arithmetic. The first column of results gives the exact rank of the matrix of binary64 numbers actually stored, which is a matrix of rational numbers; we computed it by elimination in exact rational arithmetic. The second gives the number of pivots that Gaussian elimination with partial pivoting, carried out in binary64 arithmetic, finds to be nonzero; it is the rank that \cref{alg:four-subspaces} would report if run in floating-point arithmetic with the test ``pivot $\ne0$''. The third is the result of the routine \code{linalg.matrix\_rank} of the stardust library, called with its default settings, which decides the rank by a different method, described below. The matrices are: the matrix $A$ of \cref{sec:four-subspaces}; the same matrix with the entry $a_{34}=-1$ replaced by the binary64 number $-1+2^{-52}$; the $3\times3$ matrices whose rows are $(1,2,3)$, $(4,5,6)$, $(7,8,9)$ and $(0.1,0.2,0.3)$, $(0.4,0.5,0.6)$, $(0.7,0.8,0.9)$, the latter rounded to binary64; products $XY$ of an $n\times k$ matrix $X$ and a $k\times n$ matrix $Y$ whose entries are integers drawn at random from $-9$ to $9$, which are stored exactly; and the Hilbert matrices $H_n$ of \cref{sec:pivoting}, with the entries $1/(i+j-1)$ rounded to binary64. The measurements were made with the library version and on the machine described in \cref{sec:systems-practice}.

\begin{table}
\sdfloat{%
\begin{sdtable}{colspec={llrrr}, row{1}={font=\small\bfseries}}
\toprule
matrix & size & exact rank & elimination & \code{matrix\_rank} \\
\midrule
$A$ of \cref{sec:four-subspaces} & $3\times4$ & $2$ & $2$ & $2$ \\
$A$ with $a_{34}=-1+2^{-52}$ & $3\times4$ & $3$ & $2$ & $2$ \\
entries $1,2,\ldots,9$ & $3\times3$ & $2$ & $3$ & $2$ \\
entries $0.1,0.2,\ldots,0.9$ & $3\times3$ & $3$ & $3$ & $2$ \\
$XY$, $k=3$ & $10\times10$ & $3$ & $10$ & $3$ \\
$XY$, $k=10$ & $100\times100$ & $10$ & $100$ & $10$ \\
$XY$, $k=50$ & $100\times100$ & $50$ & $100$ & $50$ \\
$H_8$ & $8\times8$ & $8$ & $8$ & $8$ \\
$H_{10}$ & $10\times10$ & $10$ & $10$ & $10$ \\
$H_{12}$ & $12\times12$ & $12$ & $12$ & $11$ \\
$H_{14}$ & $14\times14$ & $14$ & $14$ & $11$ \\
\bottomrule
\end{sdtable}}{The rank of matrices stored in binary64 arithmetic, computed in three ways: exactly, for the stored numbers; as the number of nonzero pivots of Gaussian elimination with partial pivoting in binary64 arithmetic; and by the library routine \code{linalg.matrix\_rank} with its default tolerance. The products $XY$ have random integer factors $X\in\mathbb Z^{n\times k}$ and $Y\in\mathbb Z^{k\times n}$; $H_n$ is the Hilbert matrix.\label{tab:rank-experiment}}
\end{table}

The table shows every combination of agreement and disagreement. For the matrix $A$, whose entries and multipliers are small integers and halves, the elimination is exact, and all three methods agree. Changing one entry by $2^{-52}$ raises the exact rank to $3$, in accordance with \cref{prop:rank-perturbation}\,(ii); the floating-point elimination misses the change, because after the first step the entry that should be $3+2^{-52}$ is rounded to $3$, and the last pivot, which in exact arithmetic is $2^{-52}/3$, comes out as exactly zero. For the matrix with the entries $1,\ldots,9$, which has rank $2$, the floating-point elimination finds a third pivot of $2^{-53}\approx1.1\times10^{-16}$, a rounding error that the test ``pivot $\ne0$'' cannot tell from a genuine entry. The matrix with the entries $0.1,\ldots,0.9$ is the same matrix divided by $10$, of rank $2$ as written; but the numbers $0.1,\ldots,0.9$ are not binary64 numbers (Volume~V), the stored matrix differs from the intended one by the rounding of its entries, and its exact rank is $3$. For the products $XY$, whose stored entries are exact, the floating-point elimination finds full rank in all three cases, a hundred nonzero pivots for a matrix of rank $10$: the rounding errors of the operations, not of the data, are enough. The Hilbert matrices are invertible, and so are their rounded versions, but the library reports the rank $11$ for $H_{12}$ and $H_{14}$.

The library routine is based on the singular values of the matrix, which Chapter~21 defines: numbers $\sigma_1\ge\sigma_2\ge\cdots\ge0$, as many as $\min(m,n)$, of which exactly $\operatorname{rank}A$ are nonzero in exact arithmetic, and which change by at most $\|E\|_2$ when $A$ is perturbed to $A+E$. It counts the singular values that exceed a tolerance relative to the largest one. Its answer is therefore the \term{numerical rank}: the smallest rank of a matrix within a distance of about the tolerance from the given one, which is the right question when the entries carry errors. For the diagonal matrices $\operatorname{diag}(1,s)$ the routine returned the rank $2$ for $s=4.5\times10^{-16}$ and $1$ for $s=4.4\times10^{-16}$, and for $\operatorname{diag}(1,1,s)$ it returned $3$ for $s=6.7\times10^{-16}$ and $2$ for $s=6.6\times10^{-16}$, consistent with a default tolerance of $\max(m,n)$ times the spacing $2^{-52}$ of the binary64 numbers near $1$. The numerical rank depends on the tolerance, which the user may set: with the tolerances $10^{-10}$, $10^{-12}$ and $10^{-14}$ relative to the largest singular value, the routine returned for $H_8$, $H_{10}$, $H_{12}$ and $H_{14}$ the ranks $7,8,9,9$, then $8,9,10,10$, and then $8,10,11,11$. None of these numbers is wrong; they answer different questions. The exact rank of a matrix of measured or rounded data is almost always full, by \cref{prop:rank-perturbation}\,(ii), and therefore almost never what one wants to know, while the numerical rank says how many directions the data distinguish reliably at a given level of accuracy. Chapter~21 explains the singular value decomposition on which it rests, and Volume~VIII the algorithms and the choice of the tolerance~\cite[Chap.~5]{golub2013}.

\bigskip
A vector space is a set in which vectors can be added and multiplied by scalars according to the eight rules of $\mathbb R^n$, and the examples range from columns and matrices to polynomials, functions, sequences and the solutions of linear differential equations. Its subspaces are the subsets closed under the two operations; intersections and sums of subspaces are subspaces, and the shifted subspaces, the flats, are the solution sets of consistent linear systems. A list of vectors is independent if no vector is a combination of the others, and a basis is an independent list that spans; the exchange lemma, which rests on the fact that a homogeneous system with more unknowns than equations has a nonzero solution, shows that all bases have the same number of elements, the dimension. A basis gives every vector a coordinate vector, which turns an $n$-dimensional space into a copy of $\mathbb F^n$, and coordinates with respect to two bases are related by an invertible matrix. Every matrix has four fundamental subspaces, and one elimination gives bases of all of them: the pivot columns of $A$, the nonzero rows of $R$, the special solutions and the last rows of the matrix $G$ with $GA=R$. Their dimensions $r$, $r$, $n-r$ and $m-r$ prove that a matrix and its transpose have the same rank, the rank--nullity theorem in matrix form, and the Fredholm alternative, by which a system is solvable exactly when every relation among its equations holds among its right-hand sides. The rank has many faces, from the number of pivots to the size of the largest invertible submatrix and the shortest factorisation $A=CR$, and in floating-point arithmetic it gives way to the numerical rank. Chapter~16 develops linear maps between abstract vector spaces, Chapter~18 adds orthogonality to the four subspaces, and Chapter~21 computes the numerical rank from the singular values.

\chapternotes

The idea of a vector space grew out of geometry and the theory of linear equations. Grassmann's \emph{Ausdehnungslehre} of 1844~\cite{grassmann1844} contained linear combinations, linear independence and dimension, in a language so unusual that it was hardly read. Peano gave the first axiomatic definition of a real vector space, essentially \cref{def:vector-space}, in 1888, in a book that presented Grassmann's ideas~\cite{peano1888}, and the axioms became standard only in the twentieth century, through algebra and through the function spaces of analysis. The exchange lemma is named after Steinitz, who proved it in 1913 in a paper on conditionally convergent series in several dimensions~\cite{steinitz1913}, although the idea of exchanging vectors one at a time had appeared earlier, in Grassmann's work among others~\cite{moore1995}. The term rank and the equality of row rank and column rank belong to the theory of determinants and linear equations of the nineteenth century, associated with the names of Sylvester, Kronecker and Frobenius. Moore~\cite{moore1995} and Dorier~\cite{dorier1995} trace this history. Hamel's basis of $\mathbb R$ over $\mathbb Q$ is in~\cite{hamel1905}, and the equivalence of the existence of bases with the axiom of choice is due to Blass~\cite{blass1984}. Fredholm's alternative comes from his theory of integral equations~\cite{fredholm1903}, and Kirchhoff's laws and his count of loops from his paper on electrical networks~\cite{kirchhoff1847}. The four fundamental subspaces, their dimensions and the picture of \cref{fig:four-subspaces} were made the centre of a first course by Strang, who calls the results of \cref{sec:four-subspaces} and their orthogonal complement in Chapter~18 the fundamental theorem of linear algebra~\cite{strang1993,strang2016}; the factorisation $A=CR$ is advocated by Strang and Moler~\cite{strang2022}.

Axler~\cite{axler2024} develops vector spaces, bases, dimension and quotient spaces without matrices and with complete proofs (his Chapters~1 to~3), and Halmos's classic~\cite{halmos1958} and Hoffman and Kunze~\cite{hoffman1971} are the traditional rigorous accounts. Strang~\cite[Chap.~3]{strang2016} and Meyer~\cite[Chap.~4]{meyer2000} treat the four subspaces with many examples, and Horn and Johnson~\cite[Sects.~0.1 to~0.4]{horn2013} collect the facts on rank used in matrix analysis. The incidence matrix and the Laplacian matrix of a graph are the subject of algebraic graph theory, for which Biggs~\cite{biggs1993} and Godsil and Royle~\cite{godsil2001} are standard texts, and MacWilliams and Sloane~\cite{macwilliams1977} treat linear codes over finite fields. The numerical rank and its computation by the singular value decomposition are treated by Golub and Van Loan~\cite{golub2013}.

\chapterreferences
